%!TEX root = ../thesis.tex
%*******************************************************************************
%****************************** Third Chapter **********************************
%*******************************************************************************
\chapter[Unified Framework for Conditional Sampling and Finetuning of Diffusions]{A Unified Framework for Conditional Sampling and Efficient Finetuning of Diffusions}
\label{chapter:deft}
% **************************** Define Graphics Path **************************
\ifpdf
    \graphicspath{{Chapter4/Figs/Raster/}{Chapter4/Figs/PDF/}{Chapter4/Figs/}}
\else
    \graphicspath{{Chapter4/Figs/Vector/}{Chapter4/Figs/}}
\fi

In the previous chapter, we introduced a unified framework for sampling and generative modelling based on divergences on path measures. Furthermore, our framework allowed us to propose a novel extension to AIS-based algorithms through the addition of a control term, allowing us to learn more accurate proposal and target distributions and guarantee a zero-variance importance sampling estimate.

In this chapter, we turn our attention to another important paradigm in the sampling literature: \textit{conditional generative modelling}. While denoising diffusion models have been widely successful in generating high-quality images \citep{ho2020denoising}, and designing protein backbones \citep{watson2023novo}, we are interested in the task of \textit{conditioning} the time reversal process, in order to generate samples that are subject to an observed condition. Instead of generating samples from the unconditional distribution $\pi$, we instead generate samples from the conditional distribution $\picond$, which we write either in terms of a generic reward function $r(x)$ that results in $\picond(x) \propto \hat{\pi} \, \exp (r(x))$, or an explicit likelihood model $p(y \mid x)$ for some auxiliary measurements $y$ that results in $\picond(x \mid y) \propto \hat{\pi} \, p(y \mid x)$\footnote{We might notice that this decomposition of the posterior $\picond(x \mid y)$ into a product of the prior $\pi$ and the likelihood $p(y \mid x)$ is simply Bayes' theorem.}.

Diffusion models typically approximate the time-varying score of the underlying distribution, i.e., $s_{\theta}(x, t) \approx \nabla_x \ln \pi_t(x)$. In conditional generative modelling, we instead attempt to learn a neural network $s_\theta(x, t, y)$ that approximates the score of the conditional distribution $\nabla_{x} \ln \picondt(x \mid y)$, termed the \textit{posterior score}. The time-dependent conditional distribution is defined similarly to \Cref{eq:sm_noised_distributions} as the iteratively noised version of the conditional distribution,
\[
    \picondt(x \mid y) &= \int \picond(w \mid y) \mathcal{N}(x; w, \sigma_t^2) \dd w, \\
    &\propto \pi_t(x) p(y \mid x_t). \label{eq:conditional_posterior_score}
\]
Therefore, we can write down the posterior score from \Cref{eq:conditional_posterior_score} (or using Bayes' theorem) as
\begin{align}
    \label{eq:bayes_theorem}
    \nabla_{x} \ln \picondt(x \mid y) \approx s_{\theta, t}(x) + \nabla_{x} \ln p(\m | x_t).
\end{align}
The time-dependent likelihood $\nabla_{x} \ln p(y | x_t)$ is often termed \textit{guidance} due to its interpretation to guide the reverse process to the conditioned inputs, and is unfortunately analytically intractable. To tackle this problem, several approximations for the \textit{guidance} have been proposed; see, for example, \textcite{chung2022diffusion,finzi2023user,jalal2021robust,rout2024solving,song2022pseudoinverse,song2021solving} and further discussion in Appendix~\ref{sec:related_work}. Instead of relying on the decomposition in \Cref{eq:bayes_theorem}, another line of work aims to learn the posterior score directly \citep{batzolis2021conditional,ho2022classifier}, which requires expensive training for new conditional sampling tasks, and access to large amounts of paired data points.

In the setting of conditional generation with diffusion models, our primary goal is to leverage large pre-trained foundation models which are prevalent in applications, but which typical front-end users are not able to backpropagate through, making approaches like \textcite{chung2022diffusion,finzi2023user} infeasible. This might be due to their prohibitive computation times or because they lie behind an API preventing the usage of autodiff frameworks.

In this chapter, we propose a unified framework for conditional generation using Doob's $h$-transform, a well-known result in the stochastic differential equations (SDE) literature \citep{de2021simulating,rogers2000diffusions,sarkka2019applied,ye2022first}, which we have introduced in \Cref{sec:doobs_h_transform}. We show how a large number of conditional generative modeling algorithms can be unified under the lens of our framework, as well as provide novel theoretical guarantees on the convergence of different algorithms. 

Finally, using our framework, we propose DEFT \textit{(Doob's h-transform Efficient FineTuning)}, an algorithm that estimates the time-dependent likelihood directly from data, i.e., $h^* = \nabla_{x} \ln p(\m | x_t)$, while being able to leverage an existing pre-trained unconditional model. We learn the guidance term ($h$-transform) efficiently using 1) smaller networks, and 2) a small training dataset of paired data points and corresponding observations. Furthermore, through connections to stochastic control, we propose a novel network architecture for general-purpose fine-tuning, which, in conjunction with our proposed loss, achieves competitive results across a series of inverse problems in imaging and protein design, while having a much lower computational cost.


\section{Conditioning diffusions via Doob's \texorpdfstring{$h$}{h}-transform}  \label{sec:deft_h_transform}
% NOTE: Maybe write about ergodicity a bit more clearly in CMCD section.
We have previously discussed how the Doob's $h$-transform can be used to condition the boundary points of an SDE mathematically. In this section, we explicitly show how we can perform conditioning in forward-backward SDEs that are used for generative modelling. Let us briefly recap the score-based generative modelling framework of \textcite{song2020score} and the time-reversal property of diffusion models we have discussed extensively in \Cref{chapter:cmcd} (see \Cref{def:nelson-reversal} for a refresher); we start with a forward SDE, which progressively transforms the target distribution $\pi$
\begin{align}
    \dd X_t &= f_t(X_t) \, \dd t + \sigma_t \dd W_t \quad X_0 \sim  \pi, \quad X_{0:T} \sim \ora{\P}^{\pi, f} \label{eq:forward_sde_doobs}
\end{align}
with drift $f_t$ and diffusion $\sigma_t$, until a terminal distribution $\rho_T = \ora{\P}^{\pi, f}_T$ (Recall the notation from \Cref{eq:notation_for_path_measures}, where $\rho_T$ is the marginal distribution of the forward process at time $T$). From \Cref{def:nelson-reversal}, we know that under ergodicity assumptions on the forward process so that when $T$ is large, $\rho_T$ is close to a simple distribution $\pi_T = \mathcal{N}(0, I)$. there exists a corresponding reverse SDE with corresponding drift $b_t$ \citep{anderson1982reverse}, that allows us to take samples from $\pi_T$ and denoise them to generate samples from $\pi$,
\begin{align}
    \label{eq:back_sde}
    \dd X_t &= \left( f_t(X_t) - \sigma_t^2 \nabla_{X_t} \ln p_t(X_t) \right) \,\dd t + \sigma_t \dd \ola{W}_t, \quad X_T \sim  \pi_T, \quad X_{0:T} \sim \ola{\P}^{\pi_T, b}
\end{align}
where the time flows backwards, and $b_t = f_t(X_t) - \sigma_t^2 \nabla_{X_t} \ln p_t(X_t)$.
The goal of conditional sampling is to condition the reverse SDE on a particular observation, i.e., to produce samples that satisfy constraints. For example, we might want to use \Cref{eq:back_sde} to generate samples where we already know some dimensions of the sample (e.g. knowing some pixels of the image a-priori in image inpainting). 

\subsection{Classical Noiseless Setting}

Doob's $h$-transform \citep{rogers2000diffusions,de2021simulating} provides a formal mechanism for conditioning an SDE to hit a deterministic event at a given time. We will show that existing methods for conditional generative modelling arise as approximate instances of this proposed framework. Formally, we extend the general definition of Doob's $h$-transform from \Cref{def:doobs} to the time-reversal case:

%\begin{mdframed}[style=highlightedBox]
\begin{proposition}[Doob's $h$-transform on reverse SDEs \citep{rogers2000diffusions}]
    \label{prop:htrans}
Consider the reverse SDE in \Cref{eq:back_sde}. The conditioned process $X_t | X_0 \in B$ is a solution of
\begin{align}
\dd H_t &= \left(b_t(H_t) - \sigma_t^2 \nabla_{H_t}\ln \ola{p}_{0 \mid t}(X_0 \in B | H_t)\right) \,\dd t + \sigma_t \dd \ola{W}_t, \quad H_T \sim  \pi_T \label{eq:back_sde_h},
\end{align}
with a backward drift $b_t(H_t) = f_t(H_t) - \sigma_t^2 \nabla_{H_t} \ln \ola{p}_{0 \mid t}(X_0 \in B | H_t)$, such that the backward transition density of the conditioned SDE in \Cref{eq:back_sde_h} is given by $\operatorname{Law}(H_s | H_t) = \ola{p}_{s \mid t, 0}(x_s| x_t, x_0 \in B)$ and $\mathbb{P}(X_0 \in B) = 1$. % (similarly for a forward SDE).
Furthermore, the resulting SDE will hit the specified event within a finite time $T$.
\end{proposition}
%
\begin{lightgraybox}{Remark: Notation for Conditioning}
    Here, we note that the notation for conditioning is slightly different from the notation used in \Cref{def:doobs}. In particular, we have the following notational choices:
    \begin{itemize}
        \item We use $H_t$ to denote the conditioned process, and $X_t$ to denote the unconditional process. 
        \item We refer to ``conditioned'' marginal probability distributions as $\ola{p}_{0 \mid t}(X_0 \in B | H_t)$, where we explicitly use an arrow to convey that the conditioning is being satisfied when the SDE moves in the reverse direction. This is somewhat different to the ``conditioned'' probability distribution used in \Cref{def:doobs}, where we had $p_{T \mid t}(X_T \in B | X_t)$ to denote the conditional probability distribution, without an explicit arrow, to convey that the conditioning is being satisfied when the SDE moves in the forward direction\footnote{This is in distinction to an ``unconditioned'' marginal probability distribution, which we denote as $p_{t}(X_t)$, without arrows, as the marginal probability distribution of the forward process should be equal to the marginal probability distribution of the reverse process as per Nelson's reversal in \Cref{def:nelson-reversal}.}.
        \item We denote the conditioned transition probabilities for the SDE in \Cref{eq:back_sde_h} as $\operatorname{Law}(H_s | H_t) = \ola{p}_{s \mid t, 0}(x_s| x_t, x_0 \in B), s < t$, to emphasise that the transition densities of the reverse conditioned SDE are conditioned on the boundary condition $x_0 \in B$ throughout. Similarly, we would denote the transition densities of the forward SDE as $\ora{p}_{t \mid s}(x_t| x_s), s < t$.
    \end{itemize}
\end{lightgraybox}

The function $h(t, H_t) \triangleq \ola{p}_{0 \mid t}(X_0 \in B \mid H_t)$ is referred to as the \emph{$h$-transform} \citep{rogers2000diffusions,de2021simulating}. Doob's $h$-transform is a powerful tool for conditioning an SDE, but in the formulation of \Cref{prop:htrans}, we are conditioning on a deterministic event $X_0 \in B$. Concretely, we consider events of the form $X_0 \in B$ which are described by an equality constraint $\mathcal{A}(X_0) = y$ with $\mathcal{A}$ a known \emph{measurement} (or \emph{forward}) operator and $y$ an observation, which is a common setup in inverse problems such as inpainting or super-resolution.

\begin{corollary}\label{col:inpaint}% (hard constraint)
Consider the reverse SDE in \Cref{eq:back_sde}, then it follows that
\begin{align}
\label{eq:rev_sde_htrans}
    \dd H_t &= ( b_t(H_t) - \sigma_t^2 \nabla_{H
_t}\ln \ola{p}_{0 \mid t}(\mathcal{A}(X_0) = y \mid H
_t))\,\dd t + \sigma_t \dd \ola{W}_t, % \quad X_T \sim  \dist_T
\end{align}
with $b_t(H_t) = f_t(H_t) - \sigma_t^2 \nabla_{H_t} \ln p_t(H_t)$
satisfies $\operatorname{Law}(H_s | H_t) = \operatorname{Law}(X_s | X_t, \mathcal{A} (X_0) = y)$ thus $\operatorname{Law}(H_0)  = \operatorname{Law}(X_0 | \mathcal{A} (X_0) = y)$.
\end{corollary}
%\end{mdframed}

Sampling from \Cref{eq:rev_sde_htrans} directly provides samples $x \sim \pi$ which also satisfy $\mathcal{A}(x) = y$. Crucially, this SDE is guaranteed to hit the conditioning in finite time, unlike prior equilibrium-motivated approaches \citep{chung2022diffusion,dutordoir2023neural,finzi2023user,han2022card,meng2022diffusion, song2022pseudoinverse}. 

% \paragraph{Reconstruction guidance.}
% %
% % To paraphrase \cref{col:inpaint} further, had we have access to the $h$-transform, then we would be able to solve the inpainting task.
% To get better insight into the challenge of sampling via Doob's $h$-transform in \Cref{eq:rev_sde_htrans} let us re-express the $h$-transform as
% \begin{align}
%     \ola{p}_{0 \mid t}(\mathcal{A}(X_0) = y \mid H_t) = \int \mathds{1}_{\mathcal{A}(x_0) = y} \ola{p}_{0 \mid t}(x_0 | H_t) \dd x_0,
% \end{align}
% where in the case of denoising diffusion models $\ola{p}_{0 \mid t}(x_0 | H_t)$ is the transition density of the reverse SDE \eqref{eq:back_sde}.
% In practice, one does not have access to this transition density, i.e. we can sample from this distribution, but we cannot easily get its value at a certain point. This makes it difficult to approximate the integral.
% To alleviate this, a strand of recent works \citep{chung2022diffusion,finzi2023user,rozet2023score,song2022pseudoinverse} have proposed to apply a Gaussian approximation of $\ola{p}_{0 \mid t}(x_0 | H_t) \approx \mathcal{N}(x_0 ~\vert~ \E[X_0 | X_t = H_t], \Gamma_t)$ leveraging Tweedie's formula and the pre-trained score network. This line of work is referred as \emph{reconstruction guidance}.
% We note that whilst proposing to approximate the quantity $\ola{p}_{0 \mid t}(\mathcal{A}(X_0) = y \mid H_t)$, they do not make the connection to Doob's transform and thus are unable to provide theoretical guarantees that the samples will satisfy the conditioning at finite time, such as in \Cref{col:inpaint}.
% %
% Overall, the Gaussian-based approximations of Doob's $h$-transform lead to reconstruction guidance-based approaches that simulate the following reverse SDE for individual samples at test-time:
% \begin{align}
%   \dd H_t = \left( b_t(H_t) +\sigma_t^2 \nabla_{H_t} ||y - \mathrm{A} \E[X_0 | X_t = H_t]  ||^2_{\Gamma_t}\right)\dd t + \sigma_t \dd \ola{W}_t, X_T \sim  \pi_T,  
% \end{align}
% where $\Gamma_t$ acts as a guidance scale \citep{rozet2023score,mathis2023nmd}, and $\mathrm{A}$ is a matrix if $\mathcal{A}$ is linear otherwise $\mathrm{A} = \mathrm{d} \mathcal{A}(\E[X_0 |X_t = H_t])$. The expectation $\E[X_0 |X_t = H_t]$ can be computed via Tweedie's formula utilising the already trained score network, and usually involves atleast one backpropagation step through the score network. Therefore, each step of simulating the reverse SDE often requires a backpropagation step through the score network, which is computationally expensive.

\subsection{A More General Noisy Setting}

Rather than conditioning an SDE on a deterministic event, one is often interested in a posterior arising from noisy observations (e.g. noisy inverse problems)
\begin{align}
    Y = \text{noisy}(\mathcal{A}(X_0)), \,\, X_0 \sim \pi,
\end{align}
where $\mathcal{A}$ is a forward operator, ``$\text{noisy}$'' describes a noise process and unlike the classical $h$-transform, we are not enforcing a deterministic condition such as $\mathcal{A}(X_0) = Y$. We typically assume we can evaluate and sample from the likelihood $p(y | X = x_0)$. Our goal is to sample from the posterior $p(x_0 | Y = y) = p(y | x_0) \pi(X_0) / p(y)$. Unfortunately, Doob's $h$-transform does not immediately apply to this setting, as we are not conditioning on a deterministic event. In the next section however, we extend the $h$-transform to this setting, and show how it serves as a natural framework to unify many existing methods for conditional generation.

\section{A Unifying Framework: The Generalised \texorpdfstring{$h$}{h}-transform}

Sampling from the posterior can be achieved by a novel generalisation of the $h$-transform that we propose:

\begin{proposition}[Generalised $h$-transform]
\label{prop:noisy_cond}
Given the following backwards SDE with marginals $p_t$
\begin{align} 
    \dd X_t &=  b_t(X_t) \,\dd t + \sigma_t\; \dd \ola{W}_t,  \quad X_T \sim \pi_T, \label{eq:rev_sde2trans_21}
\end{align}
then it follows that the backward SDE 
\begin{align} 
    \dd H_t = \left( b_t(H_t) - \sigma^2_t  \nabla_{H_t} \ln \ola{p}_{y \mid t}(Y=y|H_t)\right)\,\dd t + \sigma_t\; \dd \ola{W}_t, \quad H_T &\sim  \ora{\P}_T^{p(x_0 \mid y),f_t}\label{eq:rev_sde2trans_2}
\end{align}
satisfies $\operatorname{Law}(H_0) = p(x_0 | Y=y)$ with 
\[  
\ola{p}_{y \mid t}(Y=y | H_t) = \int p(Y = y | x_0) \ola{p}_{0 \mid t}(x_0 | H_t) \dd x_0, \label{eq:gen_htrans_expr}
\]
where $p(Y = y | x_0)$ is the likelihood of the noisy observation $Y = y$, and therefore \Cref{eq:gen_htrans_expr} is the $h$-transform from \Cref{prop:htrans} integrated over the likelihood $p(Y = y | x_0)$. Furthermore, the initial density for \Cref{eq:rev_sde2trans_2} is given by
\[
    H_T \sim  \ora{\P}_T^{p(x_0 \mid y),f_t} = \int \ora{p}_{T|0}(x| x_0) p(x_0|y) \dd x_0 .
\]
where $p(x_0 \mid y)$ is the posterior distribution of the data given the noisy observation $Y = y$, and $\ora{p}_{T|0}(x| x_0)$ refers to the transition density of the forward SDE with drift $f_t$. Therefore, $\ora{\P}^{p(x_0 \mid y),f_t}$ refers to the path measure for this forward SDE with drift $f_t$ but initial distribution $p(x_0 \mid y)$ (instead of the unconditional distribution $\pi(x_0)$).
%
\end{proposition}
\begin{proof}
    See Appendix \ref{app:proof_noisy_cond}.
\end{proof}

\begin{lightgraybox}{Remark: Notation for Generalised $h$-transform}
    \begin{itemize}
        \item The generalised ``conditioned'' probability distribution is given by $\ola{p}_{y \mid t}(Y=y | H_t)$, where the $y \mid t$ emphasizes that the conditioning is on the noisy likelihood $p(y | x_0)$, and the reverse arrow emphasizes that the conditioning is on the reverse direction. 
        \item The initial distribution for the reverse conditioned SDE requires some careful attention, as this is no longer equal to $\pi_T$ as we have had so far. Instead, it is given by $\ora{\P}^{p(x_0 \mid y),f_t}_T$, which refers to the terminal distribution at time $T$ of the path measure corresponding to a forward SDE specified similar to \Cref{eq:forward_sde_doobs} but with initial distribution which is given by the posterior distribution $p(x_0 \mid y)$ corresponding to the likelihood $p(y | x_0)$ (instead of the unconditional distribution $\pi$), which is given by
        \[
            \dd X_t &= f_t(X_t) \, \dd t + \sigma_t \dd W_t \quad X_0 \sim  p(x_0 \mid y), \quad X_{0:T} \sim \ora{\P}^{p(x_0 \mid y), f}.
        \]
        Therefore, for the generalised $h$-transform, we have a requirement that the conditioned reverse SDE has to have an initial distribution which corresponds to the terminal distribution of the forward SDE, if \textit{it had started from the posterior distribution $p(x_0 \mid y)$}.
    \end{itemize}
\end{lightgraybox}

In practice, the initial distribution of the reverse SDE can pose an issue, as we do not know it in closed form. However, as we will show later, we can bound the error introduced by approximating this initial distribution with something tractable, i.e. $\mathcal{N}(0, I)$, and show that it decays exponentially for increasing $T$.

To summarise \Cref{prop:noisy_cond}, the result gives a generalisation of the $h$-transform that allows sampling from posteriors; notice that it recovers the traditional $h$-transform in the no-noise setting. Whilst this more general formulation of the $h$-transform has been explored in unconditional generative modelling \citep{ye2022first}, this is the first work to cast conditional generative modelling in this light. We refer to the corresponding term $h(x_t, t, y) \triangleq \ola{p}_{y \mid t}(Y = y \mid H_t)$ as the \textit{generalised $h$-transform} henceforth.

\subsection{Existing Approaches under the Generalised \texorpdfstring{$h$}{h}-transform}

\begin{figure}[!t]
    \centering
    \includegraphics[width=\textwidth]{"/Users/shreyaspadhy/Documents/Obsidian Vault/phd-thesis/Chapter5/Figs/Vector/four_sampling_labeled.pdf"}
    \caption[Illustration of four existing approaches for conditioning within the generalised \texorpdfstring{$h$}{h}-transform framework.]{Illustration of four existing approaches for realising the boundary condition \(x_b=\hat{x}\) within the generalised \(h\)-transform framework. Grey contours show the target density \(\pi(x_a,x_b)\); the black dashed line is the constraint manifold \(x_b=\hat{x}\).  
    The \(\times\) marks are successive samples obtained during the simulation of the reverse SDE starting from an arbitrary point, and arrows indicate proposed updates.  
    \textbf{(a) Unconditional sampling} follows the reverse SDE of \(\pi\) towards the unconditional density, and ignores the constraint.  
    \textbf{(b) Replacement} simulates the same SDE but overwrites the observed coordinate after every step, \(x_b\!\leftarrow\!\hat{x}\) (shown in dotted red lines).  
    \textbf{(c) Reconstruction guidance} augments the unconditional score term in the drift (shown in blue) with Tweedie’s correction \(\sigma_t^{2}\nabla_{x_t}\ln\ola{p}_{0\mid t}(x_0\mid H_t)\) (shown in red), steering the path towards the constraint.  
    \textbf{(d) Amortised training} uses a network trained to approximate the conditional score \(\nabla_{x_t}\ln\ora{p}_{t\mid0}(x_t\mid X_0=\hat{x})\) directly, so the trajectory follows the true conditional flow from the outset.}

    \label{fig:sampling_cond}
\end{figure}

\Cref{prop:noisy_cond} provides theoretical backing to methodologies such as DPS \citep{chung2022diffusion} or $\Pi$GDM \citep{song2022pseudoinverse}, in which the reverse SDE in \Cref{eq:rev_sde2trans_2} is used to solve noisy inverse problems. While prior works have explored using Bayes' rule to decompose the conditional score, we provide rigorous arguments for intermediate steps, and carefully formalise the connection between conditional generative modelling and the $h$-transform, providing a concise result. This framework is flexible enough to also encompass prior work on conditional score matching, see e.g., \textcite{batzolis2021conditional,ho2022classifier}. In this section, we will discuss how existing approaches for conditional generation can be cast under the generalised $h$-transform framework.

\paragraph{Reconstruction Guidance.}

Let us consider the expression in \Cref{eq:gen_htrans_expr}. In practice, one does not have access to the transition density $\ola{p}_{0 \mid t}(x_0 | H_t)$ -- i.e. we can sample from this distribution, but we cannot easily get its value at a certain point. This makes it difficult to approximate the integral. \textit{Reconstruction Guidance} methods \citep{chung2022diffusion,finzi2023user,rozet2023score,song2022pseudoinverse} have proposed to approximate $\ola{p}_{0 \mid t}(x_0 | H_t)$ using Tweedie's formula, stated as follows:
\begin{definition}[Tweedie's formula]
    Let $a$, $b$ be random variables such that the conditional distribution $p(a \mid b)$ is an exponential family distribution such as $p_0(a) \exp \left(b^{\top} T(a)-\psi(b)\right)$. Then the posterior mean $\hat{b}=\mathbb{E}[b \mid a]$ should satisfy $\left(\nabla_a T(a)\right)^{\top} \hat{b}=\nabla_a \log p(a)-\nabla_a \log p_0(a)$.
\end{definition}
For DDPM models, we have that $\ora{p}_{t \mid 0}(x_t \mid x_0)$ is Gaussian (by construction of the forward process), and therefore we can apply Tweedie's formula to approximate the transition density using the posterior mean, as follows:
\[
    \hat{x}_0=\mathbb{E}\left[x_0 \mid x_t\right]=\frac{1}{\sqrt{\alpha \overline{(t)}}}\left(x_t+(1-\bar{\alpha}(t)) \underbrace{\nabla_{x_t} \log p_t\left(x_t\right)}_{\text{pre-trained } s_{\theta, t}}\right). \label{eq:tweedies_formula_ddpm}
\]
By choosing varying degrees of approximation, we can recover a large family of methods. For example, DPS \parencite{chung2022diffusion} approximates as follows:
\[
    \ola{p}_{0 \mid t}(x_0 | H_t) \approx \mathcal{N}\left(x_0 \mid \E[X_0 | X_t = H_t], \Gamma_t\right).
\]
We note that whilst proposing to approximate the quantity $\ola{p}_{y \mid t}(Y = y \mid H_t)$, these methods do not make the connection to Doob's transform and thus are unable to provide theoretical guarantees that the samples will satisfy the conditioning at finite time, such as in \Cref{col:inpaint}.

Overall, the Gaussian-based approximations of Doob's $h$-transform lead to reconstruction guidance-based approaches that simulate the following reverse SDE for individual samples at test-time:
\begin{align}
  \dd H_t = \left( b_t(H_t) +\sigma_t^2 \nabla_{H_t} ||y - \mathrm{A} \E[X_0 | X_t = H_t]  ||^2_{\Gamma_t}\right)\dd t + \sigma_t \dd \ola{W}_t, X_T \sim  \pi_T,  \label{eq:reverse_sde_reconstruction_guidance}
\end{align}
where $\Gamma_t$ acts as a guidance scale \citep{rozet2023score,mathis2023nmd}, and $\mathrm{A}$ is a matrix if $\mathcal{A}$ is linear otherwise $\mathrm{A} = \mathrm{d} \mathcal{A}(\E[X_0 |X_t = H_t])$. The expectation $\E[X_0 |X_t = H_t]$ can be computed using \Cref{eq:tweedies_formula_ddpm} utilising the already trained score network, and per reverse simulation step of the SDE, requires a backpropagation step through the score network (since we take the gradient of the expectation term in \Cref{eq:reverse_sde_reconstruction_guidance} w.r.t $X_t$, and this requires a backpropagation step through the score network).

\paragraph{Classifier-Free Guidance.}

Methodologies such as Classifier Free Guidance (CFG) \citep{ho2022classifier} do not model the forward operator explicitly. As a result, if these methods are applied to settings such as image out-painting (where the conditioning is on a subset of the random variable), these methodologies would only denoise the missing image patches. In \Cref{eq:rev_sde2trans_2}, if we expand out $b_t(H_t)$ we get
\begin{align}
    \dd H_t = \left( f_t(H_t) - \underbrace{\sigma^2 \nabla_{H_t} \ln p_t(H_t) - \sigma^2_t  \nabla_{H_t} \ln \ola{p}_{y \mid t}(Y=y|H_t)}_{\text{joint score } s_{\theta, t}(x_t, t, y)} \right)\,\dd t + \sigma_t\; \dd \ola{W}_t.
\end{align}
CFG methods learn a single model that approximates the joint score using a single score network $s_{\theta, t}(x_t, t, y)$, where occasionally, the observation $y$ is masked out, teaching the model to approximate the unconditional score $\sigma^2 \nabla_{H_t} \ln p_t(H_t)$, and when $y$ is not masked out, the model learns to approximate the joint score. At inference time, CFG methods can simulate the following reverse SDE:
\begin{align}
    \dd H_t = \left( f_t(H_t) - \left[ \gamma s_{\theta, t}(x_t, t, y) + (1-\gamma) s_{\theta, t}(x_t, t) \right] \right)\,\dd t + \sigma_t\; \dd \ola{W}_t,
\end{align}
where $\gamma$ is a scalar parameter that controls the strength of the guidance.

\paragraph{Amortised Conditional Training.}

It is also possible to learn the joint score directly using a dataset of datapoints paired with their observations $\mathcal{D}_{\text{cond}} = \{(x_0^{(i)}, y^{(i)})\}_{i=1}^N$, such as in Conditional Denoising Estimator (CDE) \citep{batzolis2021conditional}. The method learns a score network to approximate the joint score directly from the dataset $\mathcal{D}_{\text{cond}}$. In this section, we derive a more general loss function than CDE \citep{batzolis2021conditional}, where we can further amortise over the form of the measurement operator $\mathcal{A}$, which we find has even better empirical performance. 

First, we note that since $\ola{p}_{y|t}(Y = y| X_t) = \ora{p}_{t|y}(X_t |Y = y) p_0(Y = y) / p_t(X_t)$, we can re-express the Doob's transformed SDE of a reversed OU process as:
\begin{align}
    \dd H_t &= -\beta_t  \left(H_t + 2\nabla_{H_t} \ln \ora{p}_{t|0}(H_t \mid Y= y )\right)\,\dd t + \sqrt{2\beta_t }\; \dd \ola{W}_t, \;\; H_T \sim  \operatorname{Law}[X_T].
\end{align}
We then have the following proposition, which shows that the minimiser of the loss function is given by the conditional score:
\begin{proposition}[Amortised Conditional Score Matching]
    \label{prop:amortised_conditional_score_matching}
    The minimiser of
    \begin{align}
    f^*\! = \!\argmin_h\! \ \mathbb{E}_{\substack{ (X_0, Y
    ) \sim \mathcal{D}_{\text{cond}} \\ t \sim \mathrm{U}(0,T), H_t \sim \ora{p}_{t|0}(x_t|x_0)}  }\!\left[ ||h(t, H_t, y) - \nabla_{H_t}\!\ln \ora{p}_{t|0}(H_t|X_0)||^2  \!\right],%\nonumber
    \end{align}
is given by the conditional score $f^*_t(x, y) = \nabla{x} \log \ora{p}_{t|0}(x \vert Y = y)$.
\end{proposition}
\begin{proof}
    See Appendix \ref{app:proof_amortised_conditional_score_matching}.
\end{proof}
For settings where $\mathcal{A}$ varies like in image completion we sample $\mathcal{A}$ randomly and amortise it over our learned $h$-transform, i.e.  estimating $h^*_t(x, y, \mathcal{A})$. We refer to this approach as \emph{amortised conditional training}, since practically the neural network approximating the (conditional) score is amortised over $\mathcal{A}$ and $y$, instead of learning a separate network for each condition.

We note that this is different from Classifier-Free Guidance, as the joint score model is not occasionally masking out the observations $y$ in an attempt to approximate both the conditional and unconditional score, depending on the inputs to the model. Further, we $\mathcal{A}$ is assumed to be known (e.g.\ an inpainting mask).

Whilst CDE does not focus on the measurement model, they estimate a very similar quantity as our \Cref{prop:amortised_conditional_score_matching}:
\begin{align}
   f^{\mathrm{CDE}} (x, y)=  \nabla_{x} \log \ora{p}_{t|0}(x \vert Y = y ).
\end{align}
One can see that the quantities are related to one another via marginalizing the measurement model $p_{\mathcal{A}}$. This introduces a practical and conceptual difference: If we consider in/out painting as an example, the score network estimating $f^{\mathrm{CDE}}$ is not explicitly aware of where in the image the missing pixels are. As a result, it must perform inference over $\mathcal{A}$ (effectively marginalizing it) in order to know where to complete the image. This is clearly a much harder task for a single network to learn than conditioning on $\mathcal{A}$ where we provide this information.

\paragraph{Replacement Sampling.}

Replacement Sampling is a method for conditional generation that simulates the unconditional reverse SDE, but at each time step, replaces the coordinates within $X_t$ with the observations $y$. This process is only possible for settings where the forward operator $\mathcal{A}$ is a binary mask, and therefore $y$ is a partial observation of $X_0$. We consider $y = \mathcal{A}(X_0) = X_0^{[M]}$ in this case, where $X_0^{[M]}$ is the set of coordinates within $X_0$ that are observed. In Replacement Sampling, the model directly aims to sample from $p(X_t^{[\backslash M]} | X_0^{[ M]})$ and estimate this score while keeping the observed coordinates fixed. We can relate this approach to our amortised learning of Doob's $h$-transform, by noting that this is equivalent to learning the marginal conditional score
\begin{align}
    p(x_t^{[\backslash M]} | x_0^{[ M]})   = \int   \overbrace{p(x_t | x_0^{[ M]}) }^{\propto h(t,x_t)p_t(x_t) } d x_t^{[ M]}.
\end{align}
This can be viewed as estimating a marginal counterpart of our amortised $h$-transform approach. This approach can be seen in the motif scaffolding setting in RFDiffusion \citep{watson2023novo}.

\section{DEFT: Doob's \texorpdfstring{$h$}{h}-transform Efficient Finetuning}
\label{sec:training_objectives}

We have discussed prior works that either learn the posterior score from scratch, see e.g. \textcite{batzolis2021conditional,ho2022classifier}, or use approximations to the generalised $h$-transform, see e.g. \textcite{chung2022diffusion,kawar2022denoising}. These methods are either really expensive and data-hungry, or make cheap approximations that reduce the quality and fidelity of the obtained samples. Instead, in this section,we propose a method to learn the generalised $h$-transform directly from a very small dataset of paired datapoints $(x_0, y)$, while keeping the pre-trained unconditional network unchanged. We refer to this process as fine-tuning, as only the approximation to the generalised $h$-transform is learned. Our main result is given in the following theorem, where we give several representations of the generalised $h$-transform. 

\begin{theorem}[Representations of conditional SDE sampling] 
\label{theorem:representation_h}
For a given $y \sim \text{noisy}(\mathcal{A}(x_0))$, let $\ola{\Q}$ be the path measure of the conditional SDE 
\begin{align} \label{eq:vpsderev0_main}
     \dd H_t = \left( f_t(H_t) - \sigma^2_t \left(\nabla_{H_t} \ln p_t(H_t)+ h_t(H_t)\right) \right)\dd t + \sigma_t\; \dd \ola{W}_t,
\end{align}
where $H_T \sim  \ora{\P}^{p(x_0 \mid y),f_t}_T$. The generalised $h$-transform admits the following representations:
\begin{enumerate}
    \item[1)] The path measure induced by the $h$-transformed SDE satisfies $\dd\ola{\Q}^* = \dd \ola{\P} \frac{\dd p(x_0|y)}{\dd \ola{\P}_0}$,
    %\begin{align}
    %    \dd\Q^* = \dd\P \frac{\dd p(x_0|y)}{\dd \P_0} \nonumber
    %\end{align}
    where $\ola{\P}$ is the path measure of the unconditioned reverse SDE in \Cref{eq:rev_sde2trans_2} and $\ola{\P}_0$ is it's time $0$ marginal.
    \item[2)] The $h$-transform admits a \textbf{denoising score matching} representation
    \begin{align}
h_t^* &= \argmin_{h_t \in \mathcal{H}} \mathcal{L}_\text{SM}^\m(h_t) \nonumber \\ &\mathcal{L}_\text{SM}^\m(h_t) \coloneqq  \!\!\!\!\!\! \!\!
    \mathop{\mathbb{E}}_{\substack{ X_0 \sim p(x_0 | \m) \\ t\sim \mathrm{U}(0,T) \\ H_t \sim \ora{p}_{t|0}(x_t|x_0)}  }\!\!\left[ \| \left({h_t(H_t)}\!+\!\nabla_{H_t} \ln p_t(H_t) \right)\!-\!\nabla_{H_t}\ln {p}(H_t  |X_0) \|^2\right] \nonumber \end{align}
    \item[3)] The $h$-transform admits the following \textbf{stochastic control} formulation
    \begin{align}
    h_t^* =\argmin_{h_t \in \mathcal{H}} \left\{ \mathcal{L}_\text{SC}^\m(h_t) \coloneqq \E_{\ola{\Q}}\left[\frac{1}{2}\int_0^T \sigma_t^{2}|| h_t(H_t) ||^2 \dd t\right] - \E_{H_0 \sim \ola{\Q}_0} [\ln p(y| H_0)] \right\} \nonumber,
    \end{align}
    where $\ola{\Q}$ is the path measure for the conditional SDE being controlled.
    % \item[4)] The path measure induced by the $h$-tranformed SDE solves the a Schr\"{o}dinger bridge problem with boundary conditions $\ola{\Q}_0 = \ora{\P}^{p(x_0 \mid y),f_t}_T \approx \mathcal{N}(0,I)$, $\ola{\Q}_T = p(x_0|y)$ and with the unconditional process $\ola{\P}$ as its reference. 
    % \begin{align}
    %     \Q^* = \argmin_{\Q:\,\, \Q_0 = Q^{f_t}_T[ p(x_0|y)],\,\Q_T = p(x_0|y)} \KL(\Q || \P)\nonumber
    % \end{align}
    % where $\P$ as before is the path measure associated with our unconditioned reverse SDE. 
\end{enumerate}
\end{theorem}
\begin{proof}
    1) follows from \textcite{bernton2019schr,vargas2021solving}. For a proof of 2), see Appendix \ref{app:proof_representation_h_2}.
\end{proof}
Under appropriate conditions on the likelihood, the space of admissible controls $\mathcal{H}$ can be taken to be the set of $C_1$-vector fields with linear growth in space; see \textcite{nusken2021solving}. In the following sections, we will discuss the representations in 2) and 3) in more detail.  

\subsection{DEFT: Fine-tuning by score matching}
\label{sec:offline_finetuning}

The score matching objective in \Cref{theorem:representation_h}, (2), offers a simulation-free loss function to estimate the generalised $h$-transform. While the theorem's formulation focuses on learning the $h$-transform for a specific measurement $y$, this loss function can naturally be extended and amortized over the full range of measurements, i.e.,
\begin{align}
    \min_{h \in \mathcal{H}} \E_{y \sim Y}[\mathcal{L}_\text{SM}^y(h)], \label{eq:finetuning_loss_amortised}
\end{align}
to obtain $h^{*}_t(x, y) = \nabla_{x}\ln p_t(y | x)$. Further, for settings where the operator may vary, we can additionally amortise over the forward operator $\mathcal{A} \sim p$ and learn $h^{*}_t\!(x, y, \mathcal{A}) = \nabla_{x}\ln p_t(y | x, \mathcal{A})$. We exploit this to amortise over inpainting masks, see Section~\ref{sec:image_rec}, and motif scaffolding, see Section~\ref{sec:protein_design}. 
For the DDPM \citep{ho2020denoising} discretisation of the SDE and a pre-trained epsilon matching model $\epsilon_t^{\theta^*}$, the fine-tuning objective in \Cref{eq:finetuning_loss_amortised} reduces to
\begin{align}
    \label{eq:ddpm_finetuning}
    \min_\phi \mathbb{E}_{(X_0, Y), \epsilon, t} \left[ \| ( h^\phi_t(H_t, Y) + \epsilon^{\theta^*}_t(H_t)) - \epsilon \|^2 \right],
\end{align}
with $H_t = \sqrt{\Bar{\alpha}_t} X_0 + \sqrt{1 - \Bar{\alpha}_t} \epsilon$, $ (X_0, Y) \sim \mathcal{D}_{\text{cond}}, \epsilon \sim \mathcal{N}(0, I)$, where $h_t^\phi$ represents the neural network used to approximate the generalised $h$-transform. Note that the loss function in \Cref{eq:ddpm_finetuning} only requires evaluation of the pre-trained model, without needing to backpropagate through the weights $\theta^{*}$, which is often quite expensive and sometimes impossible in closed APIs. Training under the DDPM discretisation can be performed according to Algorithm~\ref{algo:finetune_dobsh_training}. Sampling with DEFT is further explained in Algorithm~\ref{algo:finetune_dobsh_sampling}, and pictorially represented in Figure~\ref{fig:condov}. As an additional insight into the behaviour of the $h$-transform that makes it more flexible and capable of modelling non-linear tasks than standard reconstruction guidance methods, we show that the $h$-transform can be interpreted as a correction term for the Tweedie estimate \citep{efron2011tweedie}.
We can express the conditional Tweedie's estimate as 
\begin{equation}
    \label{eq:h_transform_tweedie}
    \begin{split}
           \E[x_0 | x_t, y] &\approx \hat{x}_0(x_t, y) \\ &= \frac{x_t - \sqrt{1 - \bar \alpha_t} \left(h^{\phi^*}_t(x_t, y) + \epsilon^{\theta^*}_t(x_t)\right) }{\sqrt{\bar \alpha_t}} = \hat{x}_0(x_t) - \frac{\sqrt{1 - \bar \alpha_t}}{\sqrt{\bar \alpha_t}} h^{\phi^*}_t(x_t, y), 
    \end{split}
\end{equation}
where $\hat{x}_0(x_t)$ is the unconditional Tweedie estimate. \Cref{eq:h_transform_tweedie} highlights that the $h$-transform can also be interpreted as a correction factor to the unconditional denoised estimate commonly used in baselines such as DPS \parencite{chung2022diffusion}. This extension is similar to \textcite{ravula2023optimizing,zhang2023adding}. 

\begin{figure}[!t]
    \centering
    %\includegraphics[width=0.98\textwidth]{Figs/cond_overview.pdf}
    \includegraphics[width=0.95\textwidth]{"/Users/shreyaspadhy/Documents/Obsidian Vault/phd-thesis/Chapter5/Figs/Vector/DEFT_sampling_stheta_plus_network_plus_x0hat.pdf"}
    % \vspace{-0.3cm}
    \caption[DEFT reverse diffusion setup.]{DEFT reverse diffusion setup. The pre-trained unconditional diffusion model $s_t^\theta$ and the fine-tuned $h$-transform $h_t^\phi$ are combined at every sampling step. We propose a special network to parametrise the $h$-transform including the guidance term $\nabla_{\hat{x}_0} \log p(\m|\hat{x}_0)$ as part of the architecture. Here $\hat{x}_0$ denotes the unconditional denoised estimate given $s_t^\theta(x_t)$. During training, we only need to fine-tune $h_t^{\phi}$ (usually $4$-$9\%$ the size of $s_t^{\theta}$) using a small dataset of paired measurements, keeping $s_\theta^t$ fixed. During sampling, we do not need to backpropagate through either model, resulting in speed-ups during evaluation.}
    \label{fig:condov}
    % \vspace{-0.3cm}
\end{figure}

\begin{algorithm}[H]
    \caption{$\vert$ $h$-transform fine-tuning}
    \label{algo:finetune_dobsh_training}
    \begin{algorithmic}[1]
    \Require Dataset drawn from $\pi(x_0, y)$ \Comment{Dataset law $\pi$}
    \Require Noise schedule $\beta_t = \beta(t), \bar{\alpha}_t = \bar{\alpha}(t)$, parametrising process $\pi \to \pi_t$
    \Require Trained noise predictor function $\epsilon^\theta_t(x)$ with parameters $\theta$
    \Require Untrained $h$-transform $h^\phi_t(x, \hat{x}_0, y)$ with parameters $\phi$
    \Repeat
        \State $x_0 \sim \pi$
        \State $t \sim \text{Uniform}(\{1,...,T\})$
    
        % Sample measurements
        \State $y \sim p(y|x_0)$  \Comment{Simulate observations}
        % Forward noise
        \LComment{Forward noise full sample via sampling from $\ora{p}_{t \vert 0}(x_0)$}
        \State $\epsilon_t \sim \mathcal{P}_\text{noise}$
        \State $x_t \gets \sqrt{\bar\alpha_t} x_0 + \sqrt{1 - \bar\alpha_t} \epsilon_t$ 
        % Estimate noise
        \State ${\hat{\epsilon}} \gets \epsilon^\theta_t(x_t)$ \Comment{Estimate noise of sample with pretrained model}
        \State $\hat{x}_0 \gets (x_t - \sqrt{1 - \bar\alpha_t} {\hat{\epsilon}}_{\theta})/\sqrt{\bar\alpha_t}$
        \State $\hat{\epsilon}_\phi \gets h^\phi_t(x_t, \hat{x}_0, y)$ \Comment{Estimate noise of sample with $h$-transform}
        \State \text{Take gradient descent step w.r.t. $\phi$ on} 
        \Statex $\quad \nabla_{\theta}L(\epsilon, \hat{\epsilon}_\theta + \hat{\epsilon}_\phi)$
        \Comment{Typically, $L(x_\text{true}, x_\text{pred}) = || x_\text{true} - x_\text{pred} ||^2$}
    \Until converged or max epoch reached
    
    \end{algorithmic}
    \end{algorithm}
    %\subsection{Conditional Finetuning\label{app:cond_finetuning_algos}}
    %\subsection{Conditional sampling\label{app:cond_sampling_algos}}
    \begin{algorithm}[H]
    \caption{$\vert$ $h$-transform DDIM sampling}
    \label{algo:finetune_dobsh_sampling}
    \begin{algorithmic}[1]
    \Require Trained $h$-transform $h^\phi_t(x, \hat{x}_0, y)$ with parameters $\phi$
    \Require Unconditionally trained noise predictor $\epsilon^\theta_t(x_t)$ 
    \Require Noise schedule $\beta_t = \beta(t), \bar{\alpha}_t = \bar{\alpha}(t)$, parametrising process $\pi \to \pi_T$
    \Require Schedule $\sigma_t = \sigma(t)$
    \Require Observation $y$
    \LComment{Sample a starting point $x_T$}
    \State $x_T \sim \pi_T$\Comment{Often $\pi_T = \mathcal{N}(0,\mathbf{I})$}
        
    \Statex
    \LComment{Iteratively denoise for $T$ steps}
    \For{$t$ in $(T, T-1, \dots, 1)$}
            % Predict noise
            \LComment{Predict unconditional noise with learned network}
            \State ${\hat{\epsilon}} \gets \epsilon^\theta_t(x_t)$
    
            % Estimate x0hat 
            \State $\hat{x}_0 \gets \dfrac{x_t -\sqrt{1 - \bar\alpha_t} {\hat{\epsilon}}}{\sqrt{\bar\alpha_t}} $ 
    
            % Calculate h-transform 
            \State $\hat{\epsilon}_\phi \gets h^\phi_t(x_t, \hat{x}_0, \m)$
    
            \LComment{Estimate posterior noise}
            \State $\hat{\epsilon} \gets {\hat{\epsilon}} + \hat{\epsilon}_\phi$
    
            \State $\epsilon_t \sim \mathcal{N}(0, I)$ if $t>1$ else $\epsilon_t \gets 0$
            % Peform DDIM step
            \State $x_{t-1} \gets \sqrt{\bar\alpha_{t-1}} \left(\dfrac{x_t -\sqrt{1 - \bar\alpha_t} \hat{\epsilon}}{\sqrt{\bar\alpha_t}} \right) + \sqrt{1 - \bar\alpha_{t-1} - \sigma_t^2} \hat{\epsilon} + \sigma_t \epsilon_t$
        \EndFor
        \State \Return $x_0$
    \end{algorithmic}
    \end{algorithm}


\subsection{DEFT: Fine-tuning by stochastic control}
\label{sec:stochastic_control}

A limitation to the fine-tuning objective with DEFT is that it requires a small dataset of paired datapoints and measurements. In this section, we propose an alternative approach by expressing the solution to the conditional sampling problem as a stochastic optimal control objective, which is highlighted in \Cref{theorem:representation_h}, (3). This allows us to learn the $h$-transform by optimising a variational inference-type problem. Importantly, this stochastic control objective only requires the availability of a single noisy observation $y$ instead of a paired fine-tuning dataset. Further, the stochastic control objective can even be used in other conditional sampling tasks, for example in reward tilted distributions, i.e. where the goal is to sample from $\pi_{\text{cond}}(x) \propto e^{r(x)}\pi(x)$. Here $e^{r(x)}$ serves the same purpose as the likelihood, but there is no explicit measurement $y$ \citep{domingo2024adjoint}. 

In the Stochastic Control objective, one can learn the conditional score by minimising the following objective (from \Cref{theorem:representation_h}, (3)):
\begin{align}
    h_t^* =\argmin_{h_t \in \mathcal{H}} \left\{ \mathcal{L}_\text{SC}^\m(h_t) \coloneqq \E_{\ola{\Q}}\left[\frac{1}{2}\int_0^T \sigma_t^{2}|| h_t(H_t) ||^2 \dd t\right] - \E_{H_0 \sim \ola{\Q}_0} [\ln p(y| H_0)] \right\} \label{eq:stochastic_control_objective}
\end{align}
In the following proposition, we show that the minimisation of the above objective results in the correct solution to the conditional sampling problem, obtaining samples from the conditional density:
\begin{proposition}
    \label{prop:app_stochastic_control}
    The following stochastic control problem
    \begin{align}
        h^* =\argmin_{f} \E_{\ola{\Q}}\left[\frac{1}{2}\int_0^T \sigma_t^{2}|| h(H_t) ||^2 \dd t\right] - \E_{H_0 \sim \ola{\Q}_0} [\ln p(y| H_0)] \label{eq:stochastic_control_objective_2}
    \end{align}
    with 
    \begin{align} \label{eq:vpsderev0}
    &H_T \sim \ora{\P}^{p(x_0|y),f_t}_T \nonumber\\
         \dd H_t &= \left( f_t(H_t) - \sigma^2_t (\nabla_{H_t} \ln p_t(H_t)+ h_t(H_t)) \right)\,\dd t + \sigma_t\; { \dd \ola{W}}_t,
    \end{align}
    is minimised by the conditional score SDE in Equation \eqref{eq:rev_sde2trans_2}, that is
    \begin{align}
        h^*_t(x) = \nabla_{x}\ln \E_{X_0 \sim p_{0 |t}(\cdot|x)}[p(y| X_0) ] =\nabla_{x}\ln \ola{p}_{y|t}(y |x).
    \end{align}
    Furthermore, $h^*$ solves an associated half-bridge problem \citep{bernton2019schr} with the SDE in Eqn. \eqref{eq:back_sde} as its reference process and $p(x_0|y)$ as its source distribution.
\end{proposition}
\begin{proof}
    See Appendix \ref{app:proof_stochastic_control}.
\end{proof}

\begin{definition}[VP-SDE Stochastic Control and Discretisation]


Specifically, for a VP-SDE \citep{song2021solving}, the objective in \Cref{eq:stochastic_control_objective_2} reduces to
\begin{align}
    \argmin_{h} \E_{\ola{\Q}}\left[\int_0^T \frac{\beta_t}{2}|| h(\bH_t) ||^2 \dd t\right] - \E_{\bH_0 \sim \ola{\Q}_0} [\ln p(\vy| \bH_0)] \label{eq:dds}.
\end{align}
Further, we can discretise this objective easily. Let us consider the pre-trained score SDE with an added tuning network:
\begin{align} \label{eq:vpsderev2}
&\bH_T \sim \gN(0,I)\nonumber\\
    \dd \bH_t &= -\beta_t (\bH_t+ 2s_{\theta}(\bH_t) + 2h_{\phi}(\bH_t)) \,\dd t + \sqrt{2\beta_t }\;{\dd \rv{W} }_t,
\end{align}
now using an exponential-like discretisation \citep{de2022convergence}:
\begin{align} \label{eq:vpsderev2disc}
&\bH_{t_K} \sim \gN(0,I)\nonumber\\
   \bH_{t_{k-1}} &= \left(\sqrt{1- \alpha_{k}}\bH_{t_k}+ 2(1-\sqrt{1- \alpha_{k}})\left(s_{\theta^*}(\bH_{t_k}) + h_{\phi}(\bH_{t_k})\right) \right) + \sqrt{\alpha_{k}}\varepsilon_k,
\end{align}
where $\alpha_k = 1- \exp\left(\int_{t_{k-1}}^{t_k}\beta_s \dd s\right)$. We denote the distribution of the above discrete time chain as $q_\phi$. Now if we follow the sketch in Proposition 3 of \textcite{vargas2023denoising} the discretised objective then becomes:
\begin{align}
  \argmin_\phi  \E_{\bH \sim q_\phi}\left[2 \sum_{k=1}^K \frac{\lambda_k^2}{\alpha_k}|| h_{\phi}(k, \bH_{t_k})||^{2} -\ln p(\vy|\bH_{0})\right] \label{eq:objdisc}
\end{align}
where $\lambda_k = 1-\sqrt{1-\alpha_k}$. For a more stable/simple objective following \textcite{vargas2023denoising} we can make the approximation $\lambda_k = 1-\sqrt{1-\alpha_k}\approx \alpha_k / 2$ for small time steps. This leads to the following discretisation:
\begin{align} \label{eq:vpsderev3disc}
&\bH_{t_K} \sim \gN(0,I) \nonumber\\
   \bH_{t_{k-1}} &= \left(\sqrt{1- \alpha_{k}}\bH_{t_k}+ \alpha_k\left(s_{\theta^*}(\bH_{t_k}) + h_{\phi}(\bH_{t_k})\right) \right) + \sqrt{\alpha_{k}}\varepsilon_k,
\end{align}
and objective:
\begin{align}
  \argmin_\phi  \E_{\bH \sim q_\phi}\left[\sum_{k=1}^K \frac{\alpha_k}{2}|| h_{\phi}(k, \bH_{t_k})||^2 - \ln p(\vy|\bH_{0})\right]. \label{eq:stochdiscrete}
\end{align}

\end{definition}

However, the stochastic control objective is not directly applicable for high-dimensional training, as the complete chain $\{H_t\}_t$ must be kept in memory and backpropagated through or adjoint methods have to be used~\cite{li2020scalable}. In Appendix \ref{app:vargrad}, we discuss several alternatives and present experiments for scaling up the above objective, e.g., methods like VarGrad \citep{richter2020vargrad} and Trajectory Balance \citep{malkin2022trajectory}. VarGrad allows us to detach the trajectory from the gradient computation, drastically reducing the memory footprint. Further, we show initial experiments for conditional sampling in \ref{sec:stoch_cont_exp}. The stochastic control objective serves as a conceptual bridge between sampling from unnormalised densities using diffusion models \citep{vargas2023denoising,vargas2021bayesian,zhang2021path} and conditional score-based generative modelling.

\subsection{Likelihood-informed Inductive Bias}

\label{sec:network_architecture}
If the likelihood is differentiable, we can impose an inductive bias on the $h$-transform approximation. Specifically, the generalized $h$-transform can be expressed as an expectation, and we can apply the DPS approximation \cite{chung2022diffusion} as follows
\begin{align*}
    \nabla_{\vx_t} \ln \ola{p}_{y|t}(\m | \vx_t) = \nabla_{\vx_t} \ln \E_{\vx_0 \sim \ola{p}(\vx_0 | \vx_t)}[ p(\m | \vx_0)] &\approx \nabla_{\vx_t} \ln  p(\m | \E[\hat{\vx}_0|\vx_t])
    \\ &\approx  \nabla_{\vx_t} \ln p(\m | \hat{\vx}_0(\vx_t) ),
\end{align*}
where we use Tweedie's estimate from \Cref{eq:tweedies_formula_ddpm} based on the pre-trained unconditional diffusion model in the last step. The DPS approximation has been validated in many different conditional sampling tasks, so it would make for a good initialisation of the learned $h$-transform. However, the DPS approximation requires the Jacobian of the unconditional model, which is expensive to compute and known to be poorly conditioned. Further, in applications where we only have access to the forward pass of the unconditional model, the Jacobian is infeasible to compute. Similar to \textcite{poole2023dreamfusion}, we found that omitting this term still leads to an expressive architecture, while greatly reducing the computational cost. Thus, we propose the following network architecture 
\begin{align}
    \label{eq:sampling_architecture}
    h_t^\phi(\vx_t, \m) = \text{NN}_1^\phi(\vx_t, \hat{\vx}_0(\vx_t), \nabla_{\hat{\vx}_0} \ln p(\m | \hat{\vx}_0(\vx_t)), t) + \text{NN}_2^\phi(t) \nabla_{\hat{\vx}_0} \ln p(\m | \hat{\vx}_0(\vx_t)),  
\end{align}
to parametrise the $h$-transform, where the last layer of $\text{NN}_1^\phi$ is initialised with $\mathbf{0}$ and $\text{NN}_2^\phi$ is initialised to output $\mathbf{1}$. This initialisation provides a computationally efficient approximation to the $h$-transform, which still guides the sampling.

This type of network architecture has been proposed within the sampling community to apply diffusion models to normalising constant estimation \citep{phillips2024particle,vargas2023denoising,zhang2021path}. The theoretical connection to stochastic control in Section \ref{sec:stochastic_control}, motivates us further to adapt the architectures from the sampling field to the conditional generative modelling setting. We ablate the different components of our proposed architecture in Appendix~\ref{app:ablation_arch} and find that the additional components greatly improve performance empirically.

\section{Experiments}
\label{sec:exp}
We evaluate the DEFT framework from Section~\ref{sec:offline_finetuning} on both linear and non-linear natural and medical image reconstruction tasks, as well as the motif scaffolding problem in protein design. Further, in Appendix \ref{app:flowers_experiments} we provide a comparison of the conditional training framework with DEFT on the {\scshape Flowers}~\citep{nilsback08} image dataset. We provide our code \href{https://github.com/alexdenker/DEFT}{\url{https://github.com/alexdenker/DEFT}}.

\begin{figure*}
\begin{figure}[H]
\centering
\begin{subfigure}[t]{.135\textwidth}
  \includegraphics[width=\linewidth]{"/Users/shreyaspadhy/Documents/Obsidian Vault/phd-thesis/Chapter5/Figs/Raster/Figs/ImageNet/Inpainting/inp_overlay_1.png"}
\end{subfigure}%
\hfill
\begin{subfigure}[t]{.135\textwidth}
  \includegraphics[width=\linewidth]{"/Users/shreyaspadhy/Documents/Obsidian Vault/phd-thesis/Chapter5/Figs/Raster/Figs/ImageNet/Inpainting/inp_dps_1.png"}
\end{subfigure}%
\hfill
\begin{subfigure}[t]{.135\textwidth}
  \includegraphics[width=\linewidth]{"/Users/shreyaspadhy/Documents/Obsidian Vault/phd-thesis/Chapter5/Figs/Raster/Figs/ImageNet/Inpainting/inp_pigdm_1.png"}
\end{subfigure}%
\hfill
\begin{subfigure}[t]{.135\textwidth}
  \includegraphics[width=\linewidth]{"/Users/shreyaspadhy/Documents/Obsidian Vault/phd-thesis/Chapter5/Figs/Raster/Figs/ImageNet/Inpainting/inp_ddrm_1.png"}
\end{subfigure}%
\hfill
\begin{subfigure}[t]{.135\textwidth}
  \includegraphics[width=\linewidth]{"/Users/shreyaspadhy/Documents/Obsidian Vault/phd-thesis/Chapter5/Figs/Raster/Figs/ImageNet/Inpainting/inp_reddiff_paper_1.png"}
\end{subfigure}%
\hfill
\begin{subfigure}[t]{.135\textwidth}
  \includegraphics[width=\linewidth]{"/Users/shreyaspadhy/Documents/Obsidian Vault/phd-thesis/Chapter5/Figs/Raster/Figs/ImageNet/Inpainting/inp_deft_1.png"}
\end{subfigure}%

\begin{subfigure}[t]{.135\textwidth}
  \includegraphics[width=\linewidth]{"/Users/shreyaspadhy/Documents/Obsidian Vault/phd-thesis/Chapter5/Figs/Raster/Figs/ImageNet/Inpainting/inp_overlay_2.png"}
  \caption*{Ground truth}
\end{subfigure}%
\hfill
\begin{subfigure}[t]{.135\textwidth}
  \includegraphics[width=\linewidth]{"/Users/shreyaspadhy/Documents/Obsidian Vault/phd-thesis/Chapter5/Figs/Raster/Figs/ImageNet/Inpainting/inp_dps_2.png"}
  \caption*{DPS}
\end{subfigure}%
\hfill
\begin{subfigure}[t]{.135\textwidth}
  \includegraphics[width=\linewidth]{"/Users/shreyaspadhy/Documents/Obsidian Vault/phd-thesis/Chapter5/Figs/Raster/Figs/ImageNet/Inpainting/inp_pigdm_2.png"}
  \caption*{$\Pi$GDM}
\end{subfigure}%
\hfill
\begin{subfigure}[t]{.135\textwidth}
  \includegraphics[width=\linewidth]{"/Users/shreyaspadhy/Documents/Obsidian Vault/phd-thesis/Chapter5/Figs/Raster/Figs/ImageNet/Inpainting/inp_ddrm_2.png"}
  \caption*{DDRM}
\end{subfigure}%
\hfill
\begin{subfigure}[t]{.135\textwidth}
  \includegraphics[width=\linewidth]{"/Users/shreyaspadhy/Documents/Obsidian Vault/phd-thesis/Chapter5/Figs/Raster/Figs/ImageNet/Inpainting/inp_reddiff_2.png"}
  \caption*{RED-diff}
\end{subfigure}%
\hfill
\begin{subfigure}[t]{.135\textwidth}
  \includegraphics[width=\linewidth]{"/Users/shreyaspadhy/Documents/Obsidian Vault/phd-thesis/Chapter5/Figs/Raster/Figs/ImageNet/Inpainting/inp_deft_2.png"}
  \captionsetup{justification=centering}
  \caption*{DEFT}
\end{subfigure}%
\caption{Results for inpainting. We show the ground truth with the inpainting mask superimposed. 
} \label{fig:inpainting_results}
\end{figure}
\begin{table}[H]
\centering
\caption{Results on inpainting and 4x super-resolution. Best values are shown in \textbf{bold}, second best values are \underline{underlined}. We report both the total time to sample $1$k images, and the time per sample in seconds. The time to sample includes the training time for DEFT. These tasks aim to generate "natural"-looking images and therefore perceptual similarity metrics (KID, LPIPS and top-1) are more relevant. I$^2$SB (grey column) can be considered an upper bound on performance.}
\resizebox{0.81\textwidth}{!}{%
\begin{tabular}{@{}lccccc>{\columncolor{lightgray}}cccccc>{\columncolor{lightgray}}c@{}}
\toprule
& \multicolumn{6}{c}{\textbf{Inpainting}}                          & \multicolumn{6}{c}{\textbf{Super-resolution}}                    \\
 & \small DPS & \small $\Pi$GDM & \multicolumn{1}{l}{\small DDRM} & \small RED-diff & \small DEFT &  \small I$^2$SB & \small DPS & \small $\Pi$GDM & \small DDRM & \small RED-diff & \small DEFT & \small I$^2$SB \\ \midrule
PSNR $(\uparrow)$         & $21.27$ & $20.30$ & $20.72$ & $\mathbf{23.29}$   &   \underline{$22.18$} & $23.26$    & $24.83$ & $25.25$ & \underline{$25.32$} & $\mathbf{25.95}$  & $24.92$ & 23.95        \\
SSIM $(\uparrow)$         & $0.67$  & $0.82$  & $0.83$  & $\mathbf{0.87}$    &  \underline{$0.85$} & $0.86$       & $0.71$  & \underline{$0.73$}  & $0.72$  & $\mathbf{0.75}$  & $0.71$ & $0.64$         \\
KID $(\downarrow)$        & $15.28$ & $4.50$  & $2.50$  & \underline{$0.86$} &  $\mathbf{0.29}$ & $0.238$  & $10.01$ & $10.9$  & $14.0$  & \underline{$10.0$}       & $\mathbf{1.78}$ & $0.004$     \\
LPIPS $(\downarrow)$      & $0.26$  & $0.12$  & $0.14$  & \underline{$0.10$} &  $\mathbf{0.09}$ & $0.068$ & $0.16$  & \underline{$0.13$}  & $0.23$  & $0.25$       & $\mathbf{0.12}$  & $0.11$   \\
top-1 $(\uparrow)$        & $58.2$  & $67.8$  & $68.6$  & $\mathbf{72.0}$    & \underline{$71.7$} & $74.5$  &    \underline{$71.5$}  & $71.02$ & $63.9$  & $66.7$         & $\mathbf{71.9}$ & $71.6$   \\ \midrule
Time (hrs) $(\downarrow)$ & $30.72$ & \underline{$2.83$} & $\mathbf{0.33}$  & $7.86$   &  $5.2$ & N/A  & $30.72$ &   \underline{$2.83$}      &     $\mathbf{0.33}$    & $7.86$     & $5.2$ & N/A  \\ 
Time per sample (s) $(\downarrow)$ & 100.6 &  10.2       & 1.22 & 28.3 &  4.36 & N/A & 100.6 &     10.2    &    1.22     &  28.3  & 4.36 & N/A \\\bottomrule
\end{tabular}%
}
% \vspace{-0.4cm}
\label{tab:image_results_inp_sr}
\end{table}

\end{figure*}


\subsection{Image reconstruction}
\label{sec:image_rec}
We test a wide variety of both linear and non-linear image reconstruction tasks on the $256\times 256$px ImageNet dataset \cite{russakovsky2015imagenet}. We make use of a pre-trained unconditional diffusion model with $\sim500$M parameters \citep{dhariwal2021diffusion}\footnote{Checkpoints are available at \url{https://github.com/openai/guided-diffusion}}. We perform all our evaluations on a $1k$ subset of the validation set\footnote{\url{https://bit.ly/eval-pix2pix}}. For all inverse problems under consideration, the $h$-transform was trained on a separate $1$k subset of the validation set. For linear inverse problems, we compare against $\Pi$GDM \cite{song2022pseudoinverse}, DDRM \cite{kawar2022denoising}, DPS \cite{chung2022diffusion} and RED-diff \cite{mardani2023variational}. Additionally, we evaluate I$^2$SB \cite{liu20232}. The performance of I$^2$SB can be seen as an upper-bound to DEFT, as it is a conditional diffusion trained on the complete ImageNet dataset. For non-linear tasks, we only compare against DPS and RED-diff as both $\Pi$GDM and DDRM are not directly applicable to non-linear forward operators. For DEFT we make use of the DDIM sampling scheme with $100$ time steps \cite{song2021denoising}. For the comparison methods we used the same hyperparameters as in \cite{mardani2023variational} without further tuning, including the number of sampling steps (1000 for DPS and RED-Diff, 20 for DDRM and 100 for $\Pi$GDM).

We compute PSNR and SSIM, which are commonly used distortion measures, along with perceptual metrics such as Learned Perceptual Image Patch Similarity (LPIPS) \cite{zhang2018unreasonable}, Kernel Inception Distance (KID) \cite{binkowski2018demystifying}, and top-1 classifier accuracy of a pre-trained ResNet50 model \cite{he2016deep}. There is a well-known tradeoff between optimising distortion metrics versus perceptual quality, and depending on the task, one may wish for better performance along one axis at the cost of the other. For natural image tasks involving in-painting and super-resolution, it is common to prefer "natural"-looking images, which score better on perceptual similarity, whereas for tasks involving (medical) image reconstruction
%of an operator, such as HDR, phase retrieval, non-linear deblurring and medical image reconstruction, 
it is standard to prefer a lower distortion metric \cite{blau2018perception}. Further, we calculate the total time (including training for DEFT) for evaluation $1$k validation images in the "Time (hrs)" row. Furthermore, we report the effective time taken to sample a single image in the "Time per sample (s)" row. This time is calculated by fitting the largest batch size of validation images that fit on a single A100 GPU and dividing the time taken for the batch by the batch size. 

\begin{figure*}   
\begin{figure}[H]
\centering
\begin{subfigure}[t]{.162\textwidth}
  \includegraphics[width=\linewidth]{"/Users/shreyaspadhy/Documents/Obsidian Vault/phd-thesis/Chapter5/Figs/Raster/Figs/ImageNet/NonLinearDeblurring/blur_gt_1.png"}
\end{subfigure}%
\hfill
\begin{subfigure}[t]{.162\textwidth}
  \includegraphics[width=\linewidth]{"/Users/shreyaspadhy/Documents/Obsidian Vault/phd-thesis/Chapter5/Figs/Raster/Figs/ImageNet/NonLinearDeblurring/blur_deg_1.png"}
\end{subfigure}%
\hfill
\begin{subfigure}[t]{.162\textwidth}
  \includegraphics[width=\linewidth]{"/Users/shreyaspadhy/Documents/Obsidian Vault/phd-thesis/Chapter5/Figs/Raster/Figs/ImageNet/NonLinearDeblurring/blur_dps_1.png"}
\end{subfigure}%
\hfill
\begin{subfigure}[t]{.162\textwidth}
  \includegraphics[width=\linewidth]{"/Users/shreyaspadhy/Documents/Obsidian Vault/phd-thesis/Chapter5/Figs/Raster/Figs/ImageNet/NonLinearDeblurring/blur_reddiff_1.png"}
\end{subfigure}%
\hfill
\begin{subfigure}[t]{.162\textwidth}
  \includegraphics[width=\linewidth]{"/Users/shreyaspadhy/Documents/Obsidian Vault/phd-thesis/Chapter5/Figs/Raster/Figs/ImageNet/NonLinearDeblurring/blur_deft_1.png"}
\end{subfigure}%
 
\begin{subfigure}[t]{.162\textwidth}
  \includegraphics[width=\linewidth]{"/Users/shreyaspadhy/Documents/Obsidian Vault/phd-thesis/Chapter5/Figs/Raster/Figs/ImageNet/NonLinearDeblurring/blur_gt_2.png"}
  \caption*{Ground truth}
\end{subfigure}%
\hfill
\begin{subfigure}[t]{.162\textwidth}
  \includegraphics[width=\linewidth]{"/Users/shreyaspadhy/Documents/Obsidian Vault/phd-thesis/Chapter5/Figs/Raster/Figs/ImageNet/NonLinearDeblurring/blur_deg_2.png"}
  \caption*{Measurements}
\end{subfigure}%
\hfill
\begin{subfigure}[t]{.162\textwidth}
  \includegraphics[width=\linewidth]{"/Users/shreyaspadhy/Documents/Obsidian Vault/phd-thesis/Chapter5/Figs/Raster/Figs/ImageNet/NonLinearDeblurring/blur_dps_2.png"}
  \caption*{DPS}
\end{subfigure}%
\hfill
\begin{subfigure}[t]{.162\textwidth}
  \includegraphics[width=\linewidth]{"/Users/shreyaspadhy/Documents/Obsidian Vault/phd-thesis/Chapter5/Figs/Raster/Figs/ImageNet/NonLinearDeblurring/blur_reddiff_2.png"}
  \caption*{RED-diff}
\end{subfigure}%
\hfill
\begin{subfigure}[t]{.162\textwidth}
  \includegraphics[width=\linewidth]{"/Users/shreyaspadhy/Documents/Obsidian Vault/phd-thesis/Chapter5/Figs/Raster/Figs/ImageNet/NonLinearDeblurring/blur_deft_2.png"}
  \captionsetup{justification=centering}
  \caption*{DEFT}
\end{subfigure}%
\caption{Results for non-linear deblurring. We show both the ground truth, the measurements and samples for DPS, RED-diff and DEFT. DEFT is able to reconstruct high-quality images.}  \label{fig:non-linear_deblur}
\end{figure}
\begin{table}[H]
\centering
\caption{Results on different non-linear image reconstruction tasks. Best values are shown in \textbf{bold}, second best values are \underline{underlined}.}
\resizebox{0.81\textwidth}{!}{%
\begin{tabular}{lccccccccc}
\toprule
 &
  \multicolumn{3}{c}{\textbf{HDR}} &
  \multicolumn{3}{c}{\textbf{Phase retrieval}} &
  \multicolumn{3}{c}{\textbf{Non-linear Deblurring}} \\
 &
  \multicolumn{1}{l}{\small DPS} &
  \small RED-diff &
  \small DEFT &
  \multicolumn{1}{l}{\small DPS} &
  \small RED-diff &
  \small DEFT &
  \multicolumn{1}{l}{\small DPS} &
  \small RED-diff &
  \small DEFT \\ \midrule
PSNR $(\uparrow)$ &
  $7.94$ &
  \underline{$25.23$} &
  $\mathbf{28.51}$ &
  $9.99$ &
  \underline{$10.53$} &
  $\mathbf{13.03}$ &
  $17.57$ &
  \underline{$21.21$} &
  $\mathbf{25.16}$ \\
SSIM $(\uparrow)$ &
  $0.21$ &
  \underline{$0.79$} &
  $\mathbf{0.89}$ &
  $0.12$ &
  \underline{$0.17$} &
  $\mathbf{0.32}$ &
  $0.39$ &
  \underline{$0.53$} &
  $\mathbf{0.79}$ \\
KID $(\downarrow)$ &
  $272.5$ &
  \underline{$1.2$} &
  $\mathbf{0.10}$ &
  \underline{$93.2$} &
  $114.0$ &
  $\mathbf{80.89}$ &
  \underline{$12.89$} &
  $66.8$ &
  $\mathbf{0.34}$ \\
LPIPS $(\downarrow)$ &
  $0.72$ &
  \underline{$0.1$} &
  $\mathbf{0.04}$ &
   $0.66$ &
  \underline{$0.6$} &
  $\mathbf{0.52}$ &
  \underline{$0.42$} &
  \underline{$0.42$} &
  $\mathbf{0.09}$ \\
top-1 $(\uparrow)$ &
  $4.0$ &
  \underline{$68.5$} &
  $\mathbf{74.0}$ &
   $1.5$ &
  $\underline{7.2}$ &
  $\mathbf{13.1}$ &
  \underline{$30.2$} &
  $23.5$ &
  $\mathbf{69.9}$ \\ \midrule
Time (hrs) $(\downarrow)$ &
  $30.7$ &
  \underline{$7.9$} &
  $\mathbf{5.2}$ &
   $30.7$ &
  \underline{$7.9$} &
  $\mathbf{5.2}$ &
  $30.9$ &
  \underline{$8.1$} &
  $\mathbf{5.2}$ \\ 
Time per sample (s) $(\downarrow)$ &
  $100.4$ &
  \underline{$28.3$} &
  $\mathbf{4.4}$ &
   $100.6$ &
  \underline{$28.4$} &
  $\mathbf{4.4}$ &
  $101.2$ &
  \underline{$30.4$} &
  $\mathbf{4.6}$ \\
  
  \bottomrule 
\end{tabular}%
}
% \vspace{-0.3cm}
\label{tab:image_results_nonlinear}
\end{table}
\end{figure*}

\paragraph*{Inpainting} First, we evaluate DEFT on the linear inverse problem of image inpainting. We make use of the inpainting masks for the 1k subset used by \cite{saharia2022palette}\footnotemark[2], which includes masks that obscure $20\!\%-\!30\%$ of the image. Results are shown in Table~\ref{tab:image_results_inp_sr}, including the computational time for sampling all $1000$ validation images. For DEFT, this computational time additionally includes the $3.9$ hrs of training time of the $h$-transform additionally with the $1.2$ hrs of evaluation. Even with the added training time, we reduce the overall computational time for DEFT, compared to DPS and RED-diff. A visual comparison is provided in Figure~\ref{fig:inpainting_results}. Further, in Figure~\ref{fig:inp_diversity} in the Appendix, we show the diversity of samples using different initial seeds. Even though $\Pi$GDM and DDRM are faster methods, they perform significantly worse, and are only applicable for linear inverse problems. Inpainting is a task that prefers "natural"-looking image samples, and DEFT outperforms all other methods on perceptual metrics such as LPIPS and KID, being a close second on top-1 accuracy. 

\paragraph{Super-resolution} For another linear inverse problem, we evaluate $4$x noiseless super-resolution. Here, the forward operator is given by a bicubic downsampling. The results are presented in Table~\ref{tab:image_results_inp_sr}. While DEFT has a lower PSNR compared to the baseline methods, we see significant improvement on perceptual quality metrics (KID, LPIPS, and top-1 accuracy). We show visual results in Figure~\ref{fig:sr_results}.

\paragraph{High dynamic range} For the first non-linear tasks, we make use of the high dynamic range (HDR) task described in \cite{mardani2023variational}. Here, the forward operator is given by $\Pm(\vx) = \text{clip}(2\vx; -1, 1)$, where $\vx$ denotes the RGB image scaled to the range $[-1,1]$. The results are presented in Table~\ref{tab:image_results_nonlinear}. We observe that DPS struggles with this specific non-linear tasks, while DEFT achieves good results. We show a visual comparison in the Appendix, see Figure~\ref{fig:hdr}.

\paragraph{Phase retrieval} The goal in phase retrieval is to recover the image from intensity measurements only, i.e., the forward operator is given by $\Pm(\vx) = | \mathcal{F} \vx |$, with $\mathcal{F}$ as the Fourier transform. We study the same 2x oversampling setting as in \cite{chung2022diffusion,mardani2023variational}. Phase retrieval is a challenging non-linear inverse problem, as the forward operator is invariant to translations, global phase shifts and complex conjugation. In Figure~\ref{fig:phase_retrieval} in the appendix, we show samples for different initial seeds and observe a wide variety of image quality. We also observe this behaviour for the baseline methods. However, DEFT is able to achieve better performance compared to RED-diff and DPS, see also Table~\ref{tab:image_results_nonlinear}. However, there is room for further improvement to achieve good reconstructions on a consistent basis.

\paragraph{Non-linear deblurring} The non-linear deblurring task was originally proposed by \cite{chung2022diffusion}. Here, the forward operator is defined by a trained neural network \cite{tran2021explore}, resulting in a highly non-linear blur. Quantitative results are presented in Table~\ref{tab:image_results_nonlinear}. This non-linear reconstruction task was also evaluated for RED-diff in \cite{mardani2023variational}. However, we found that the forward operator of the original implementation\footnote{\url{https://github.com/NVlabs/RED-diff/tree/master}} leads to a nearly trivial reconstruction task. In Appendix~\ref{app:blur_mardani}, we show results with the code from \cite{mardani2023variational}, while Table~\ref{tab:image_results_nonlinear} shows the results with our implementation of the forward operator. Further, in Figure~\ref{fig:non-linear_deblur} we provide a visual comparison, where DEFT is able to recover the ground truth quite faithfully.

\paragraph{Ablation: Size of fine-tuning dataset}
As DEFT requires a dataset for fine-tuning, we ablate the number of training samples. We trained DEFT on a subset of 10, 100 and 200 ImageNet images for Inpainting. We see improvements of all metrics, when training on a larger dataset. The results are presented in Table~\ref{tab:inpainting_number_of_samples}. For the KID, we outperform RED-diff (KID: 0.86) even when trained on only 200 images. However, even with 10 images, we perform quite competitively, showcasing that our method is very sample-efficient when it comes to learning a conditional transform.
%\vspace{-0.5cm}

\begin{table}[th]
\centering
\caption{Varying the size of the fine-tuning dataset for DEFT for Inpainting on ImageNet.}
\resizebox{0.4\textwidth}{!}{%
\begin{tabular}{lcccc}
\toprule
 \multicolumn{5}{c}{\textbf{DEFT on ImageNet for Inpainting}} \\ 
\small Number of images & \small 10 & \small 100 & \small 200 & \small 1000 \\ \midrule
PSNR $(\uparrow)$ & $20.87$ & $20.99$ & $22.11$ & $22.18$  \\
SSIM $(\uparrow)$ & $0.83$ & $0.84$& $0.847$ & $0.85$  \\
KID $(\downarrow)$  & $1.85$ & $0.978$ & $0.401$ & $0.29$  \\
LPIPS $(\downarrow)$  & $0.123$ & $0.112$ & $0.096$ & $0.09$   \\
top-1 $(\uparrow)$  & $68.8$ & $69.6$ & $70.6$ & $71.7$  \\
  \bottomrule 
\end{tabular}%
}%
% \vspace{-0.5cm}
\label{tab:inpainting_number_of_samples}
\end{table}




% \newpage
% \subsection{Computed tomography}
% \label{sec:ct}

% \begin{table}[H]
% \centering
% \caption{Results for CT on \textsc{AAPM} and \textsc{LoDoPab-CT} and sampling time per image on a single GeForce RTX 3090. Best values are shown in \textbf{bold}, second best values are \underline{underlined}. For DEFT we use $100$ DDIM steps, while RED-diff and DPS use $1000$ time steps.}
% \resizebox{0.5\textwidth}{!}{%
% \begin{tabular}{lcccccc}
% \toprule
%      & \multicolumn{3}{c}{AAPM} & \multicolumn{3}{c}{LoDoPab-CT} \\
%      & DPS  & RED-diff & DEFT & DPS &  RED-diff & DEFT \\ \midrule
%  PSNR & $33.11$ & $\mathbf{34.85}$ & \underline{$34.73$} & $34.16$ & \underline{$34.95$} & $\mathbf{35.81}$\\
%   SSIM & \underline{$0.885$} & $0.865$ & $\mathbf{0.887}$ & $0.846$ & \underline{$0.849$} & $\mathbf{0.876}$ \\
%   Time (s) & $208.9$ & \underline{$83.4$} & $\mathbf{16.3}$ & $156.8$ & \underline{$70.1$} & $\mathbf{13.8}$ \\ 
% \bottomrule 
% \end{tabular}%
% }
% %\vspace{-1cm}
% \label{tab:ct_results}
% \end{table}


% We are evaluating DEFT both on the 2016 American Association of Physicists in Medicine
% (AAPM) grand challenge dataset \citep{mccollough2017low}, and the LoDoPab-CT dataset \citep{leuschner2021lodopab}, for details see Appendix~\ref{app:experimental_details_images}. 
% For the unconditional models we make use of the attention U-Net architecture \cite{dhariwal2021diffusion}. For the model trained on AAPM, we use exactly the same architecture ($\approx374$ params.) as in \cite{chung2022improving}, while for LoDoPab-CT we use a smaller model ($\approx133$M params.). For the forward operator, we use a parallel-beam radon transform with $60$ angles and add Gaussian noise with $\sigma_y = 1.0$, which corresponds to approx. $3.5\%$ relative noise. We compare against DPS \citep{chung2022diffusion} and RED-diff \citep{mardani2023variational}, where the parameters were obtained using a grid search on a subset of the validation set to maximise the PSNR. In Table~\ref{tab:ct_results} we present PSNR and SSIM, in addition to the sampling time, and provide a visual comparison Figure~\ref{fig:ct_lodopab}.  
% %We focus on these two quality measures, as for medical imaging minimising the distortion is more important than the perceptual quality \cite{blau2018perception}.
% For both datasets, we choose the same DEFT architecture with about $23$M parameters.%, with results in $6\%$ and $17\%$ the size of the unconditional model for AAPM and LoDoPab-CT, respectively. %For LoDoPab-CT, we $h$-transform was trained on a set of $3522$ images, while for AAPM we trained the $h$-transform on the same dataset as the unconditional model. 
%  In the Appendix~\ref{sec:add_res}, we perform an ablation regarding the parametrisation of DEFT, see Table~\ref{tab:ct_ablation}. In particular, these results show the necessity of providing the unconditional Tweedie estimate $\hat{\vx}_0$ as input to the $h$-transform in~\eqref{eq:sampling_architecture}. We observe almost a $8$dB difference in PSNR for models without the Tweedie estimate and the log-likelihood term.
 
% \begin{figure*}
% \begin{figure}[H]
% \centering
% \begin{subfigure}[t]{.2\textwidth}
%   \includegraphics[width=\linewidth]{"/Users/shreyaspadhy/Documents/Obsidian Vault/phd-thesis/Chapter5/Figs/Raster/Figs/CT/LoDoPab/gt_100.png"}
%   \caption*{Ground truth}
% \end{subfigure}%
% \hfill
% \begin{subfigure}[t]{.2\textwidth}
%   \includegraphics[width=\linewidth]{"/Users/shreyaspadhy/Documents/Obsidian Vault/phd-thesis/Chapter5/Figs/Raster/Figs/CT/LoDoPab/x_dps_100.png"}
%   \caption*{DPS}
% \end{subfigure}%
% \hfill
% \begin{subfigure}[t]{.2\textwidth}
%   \includegraphics[width=\linewidth]{"/Users/shreyaspadhy/Documents/Obsidian Vault/phd-thesis/Chapter5/Figs/Raster/Figs/CT/LoDoPab/x_reddiff_100.png"}
%   \caption*{RED-diff}
% \end{subfigure}%
% \hfill
% \begin{subfigure}[t]{.2\textwidth}
%   \includegraphics[width=\linewidth]{"/Users/shreyaspadhy/Documents/Obsidian Vault/phd-thesis/Chapter5/Figs/Raster/Figs/CT/LoDoPab/x_h_100.png"}
%   \caption*{DEFT}
% \end{subfigure}%
% \caption{Reconstructions for computed tomography on LoDoPab-CT} \label{fig:ct_lodopab}
% \end{figure}

% \end{figure*}






\subsection{Conditional protein design: motif scaffolding}\label{sec:protein_design} 
We evaluate DEFT on the contiguous motifs of the RFDiffusion benchmark \cite{watson2023novo}. In this motif scaffolding task, we sample protein $C_\alpha$ atom coordinates $\mx \in \R^{d}$ such that the generated backbone contains a targeted motif, i.e. a subset of $C_\alpha$ coordinates $\my \in \R^{n}$, similar to an image outpainting task. The forward operator is therefore given by $\my=\Pm(\mx) = \mathrm{A} \mx$, where $\mathrm{A} \in \{0,1\}^{n \times d}$ denotes a masking matrix which only selects the $n$ observed $C_\alpha$ coordinates.

%As a basis for our motif scaffolding experiments, 
We leverage the pretrained Genie diffusion model which is an unconditional model for protein backbone generation \cite{lin2023generating}. To apply DEFT to it, we use a downsized version of the unconditional base model as our $h$-transform model which only uses 200k instead of the original 4.1M parameters. To adopt this model to the DEFT algorithm, we modify the SE(3)-invariant encoder by adding additional conditional pair feature networks for the motif coordinates as well as the unconditional Tweedie estimate $\hat{\vx}_0$, similar to the setting in the previous experiments. As per Section~\ref{sec:network_architecture}, we add a time-dependent likelihood approximation term to the $h$-transform network. We train the $h$-transform network on the same SCOPe dataset as in \cite{didi2023framework}. More details on the training details can be found in App.~\ref{app:protein_motif_scaffolding_details}. We compare DEFT against DPS \cite{chung2022diffusion} and a previously published version of Genie that was trained in an amortised fashion \cite{didi2023framework}. The guidance parameter of DPS was tuned over 5 different experiment runs. The amortised model serves here as an upper limit of how well DEFT can perform with Genie as a base model.

%\begin{wrapfigure}{r}{0.4\textwidth}
%    \centering
%    \includegraphics{Figs/proteins.pdf}
%    \caption{Conditional protein designs in yellow with target motif 3IXT in blue.}
%    \label{fig:enter-label}
%\end{wrapfigure}

The overall in-silico success, defined by scRMSD $< 2$\AA~ and motifRMSD $< 1$\AA, is provided in Figure~\ref{fig:rfdiff_comp}. In the Appendix, we provide a detailed breakdown of these results, see Figure~\ref{fig:deft_9}. Further, in Figure~\ref{fig:protein_compare} and Figure~\ref{fig:protein_compare_samples} we provide a comparison of the task 1YCR for the different methods. We observe that DEFT outperforms DPS, solving 10 out of the 12 tasks compared to only 5 tasks for DPS. While it has lower success rates than the amortised model, it still solves all but two tasks in that benchmark with only 9\% of the parameter count and significantly shorter training time compared to the amortised model ($800$ epochs for DEFT vs $2100$ epochs for amortised). The low performance of DPS indicates that the base Genie model is limiting the performance here and may partly explain the performance difference between DEFT and amortised training. Exploring DEFT with a more capable base model is therefore another promising avenue for research. Excitingly, the lower training time and data requirements of DEFT enable fine-tuning a model on specific protein families for particular applications, a task that is left for future work. 

%The task of motif scaffolding in our protein setting amounts to sampling protein C alpha atom coordinates $\mx \in \R^{d}$ such that it contains a given subset of C alpha coordinates $\my \in \R^{n}$, i.e.\ $\my=\Pm(\mx) = \mathrm{A} \mx$, where $\Pm \in \{0,1\}^{n \times d}$ is a masking matrix which selects $n$ observed C alpha coordinates.
%We perform two sets of motif scaffolding experiments. We firstly compare the conditional training approach, i.e., {\scshape amortised} training in Appendix~\ref{app:amort_cond_training},  to {\scshape replacement} and {\scshape reconstruction guidance}. Upon observing that {\scshape amortised} performs significantly better, we then dive into a more detailed analysis of this method on the RFDiffusion benchmark, as well as a new SCOPe-based benchmark that is created from a hierarchical structure and sequence-based split described below. 

\vspace{-0.2cm}
\begin{figure}[tbh]
\centering
\hspace*{-1.3cm}
\begin{minipage}[b]{0.55\textwidth}
\includegraphics[width=0.98\textwidth]{"/Users/shreyaspadhy/Documents/Obsidian Vault/phd-thesis/Chapter5/Figs/Raster/Figs/insilicio_success.pdf"}
% \put(-80,10){\rotatebox{90}{hello}}
% \caption{Samples conditioned on inner patch.}
\vspace{-0.3cm}
\caption{Comparison of DPS, DEFT and amortised training for motif scaffolding for 12 contiguous targets. 4\% and 9\% are the relative sizes of the h-transform compared to the unconditional model.\label{fig:rfdiff_comp}}
\end{minipage}
% \hfill
\hspace*{0.1cm}
\begin{minipage}[b]{0.35\textwidth}

\adjustbox{max width=1.26\textwidth}{
\centering
% \small
% \vspace{-0}
\begin{tabular}{lccc}
\toprule
\scshape Metric & \scshape DPS & \scshape DEFT & \scshape Amortised \\
\midrule
%\textbf{\scshape SCOPe} \\
%\% Designable $(\uparrow)$ & $0.34_{\pm 0.01}$ & $0.27_{\pm 0.01}$ & $0.28_{\pm 0.01}$\\
%LPIPS $(\downarrow)$ & $0.25_{\pm 0.00}$ & $0.29_{\pm 0.01}$ & $0.33_{\pm 0.01}$\\
%\midrule
\% Success $(\uparrow)$ & $1.3$ & $9.2$ & $24.5$\\
\% scRMSD < 2 \AA $(\uparrow)$ & $44.3$ & $28.9$ & $42.$\\
\% mRMSD < 1 \AA $(\uparrow)$ & $4.1$ & $24.0$ & $45.8$\\

\bottomrule
\end{tabular}
}
\vspace{0.45cm}
% \vspace{0.05cm}
\vspace{-0.3cm}
\captionof{table}{{\scshape RFDiff} benchmark metrics (averaged over the 11 targets, 100 samples each). Success: scRMSD < 2\AA, motifRMSD < 1\AA. Details in \cref{sec:protein_design}.}
\label{tab:image-outpainting_prot}
\end{minipage}
\vspace{-0.4cm}
\end{figure}
\label{fig:protein_comp}

%\paragraph{Data}\label{app:scope_data}
%We evaluate on the RFDiffusion motif benchmark \citep{watson2023novo} and on a self-curated SCOPe benchmark based on a hierarchical structure-based split, jointly with a sequence-similarity-based split.
%For the RFDiffusion benchmark, we tested all sequence-contiguous motifs, resulting in 11 different motif design tasks.
%Our method readily extends to the non-contiguous motif setting and future work will address this in more detail.
%The performance on each of these targets is depicted in~\cref{fig:rfdiff_comp}.

%\begin{figure}[htb!]
%    \centering
%    \includegraphics[width=0.9\textwidth]{Figs/insilicio_success.pdf}\vspace{-0.3cm}
%    \caption{Comparison of DPS, DEFT and amortised training for motif scaffolding for 12 contiguous targets. For the motifs marked with *, we had to shorten the sampled scaffold ranges on both sides of the motif from 0-65 (0-63 for TMRX80) to 0-50 since we trained our version of Genie only on protein generation up to a length of 128 residues.}
%    \label{fig:rfdiff_comp}
%\end{figure}

%\paragraph{Results}
%We evaluate all three approaches on the continuous motifs from the {\scshape RFDiff}usion motif benchmark \citep{watson2023novo}. For the {\scshape Amortised} approach we retrain the Genie model \citep{lin2023generating} in an amortised fashion (Alg.~{\ref{algo:cond_dobsh}}), while for the {\scshape R. Guidance} and {\scshape Replacement} methods we used the publicly available unconditional model.
%We observe that amortised training outperforms the other approaches, especially replacement sampling (Fig.~\ref{fig:protein_comp}).


% and on a self-curated {\scshape SCOPe} benchmark based on a hierarchical structure-based split (App.~\ref{app:scope_data}).
%

% \begin{figure}
%     \centering
%     \includegraphics[width=0.5\textwidth]{Figs/6e6r_motifs.pdf}
%     \caption{Image Generation via Doob's $h$-transform sampling and training.}
%     \label{fig:image_generation}
% \end{figure}

% \section{Conditional Protein Design: Motif Scaffolding}
% \label{protein_design}

% \begin{figure}
%     \centering
%     \includegraphics[width=0.5\textwidth]{Figs/6e6r_motifs.pdf}
%     \caption{Protein Motif Scaffoldin via Doob's $h$-transform sampling and training.}
%     \label{fig:protein_design}
% \end{figure}


% \section{Discussion}
% \label{discussion}

\subsection{Scaling up the Control Objective} \label{app:vargrad}
Naively trying to minimise Eqn.~\eqref{eq:stochdiscrete} is demanding, as the full chain has to be kept in memory, which is infeasible for high dimensional problems. To alleviate this problem, one could make use of the stochastic adjoint sensitivity method \citep{li2020scalable}, in which an adjoint SDE is solved to estimate the gradients of the stochastic control loss in Theorem~\ref{theorem:representation_h} 3). This method has the advantage of a constant memory cost. However, the computational cost increases as both the reverse SDE and the adjoint SDE must be simulated. Instead, we discuss two alternative approaches to reduce the memory requirements.

\paragraph{VarGrad} We can make use of a VarGrad \citep{richter2020vargrad, richter2023improved} type loss to reduce the memory requirements. In contrast to the KL loss of Eqn.~\eqref{eq:stochastic_control_objective_2}, then the VarGrad loss is given by:
\begin{align}
\label{eq:var_grad_stochastic_control}
  D_{\mathrm{logvar}}(\mathbb{Q}, \mathbb{P} ; \mathbb{W}) = \mathbb{E}_{\mH^{g_t}_{0:T} \sim \mathbb{W}} \left[\left(\ln   \frac{\dd  \mathbb{Q} }{\dd  \mathbb{P}}(\mH^{g_t}_{0:T})  - \mathbb{E}\left[\ln   \frac{\dd  \mathbb{Q} }{\dd  \mathbb{P}}(\mH^{g_t}_{0:T})\right] \right)^2\right], 
\end{align}
where $\mathbb{Q}$ and $\mathbb{P}$ are defined as in the proof of Proposition~\ref{prop:app_stochastic_control}, i.e., $\mathbb{Q}$ is given by the conditional SDE and $\mathbb{P}$ is given by the unconditional SDE. The RND in Eqn.~\eqref{eq:rnd_vargrad} is evaluated at the trajectory of a reference process $\mathbb{W} = \mathrm{Law}(\mH^{g_t}_{0:T})$, given by
\begin{align} 
\bH_T &\sim Q_T^{f_t}[p(\vx_0|\vy)] \nonumber\\
     \dd \bH_t &= \left( f_t(\bH_t) - \sigma^2_t (\nabla_{\bH_t} \ln p_t(\bH_t)+ g_t(\bH_t)) \right)\,\dd t + \sigma_t\; \bwd{ \dd \rv{W}}_t.
\end{align}
The Radon–Nikodym derivative (RND) in Eqn~\eqref{eq:var_grad_stochastic_control} can be evaluated as (Using the RND for time reverse SDEs see Equation 64 in \citep{vargas2024transport}):
\begin{align}
\label{eq:rnd_vargrad}
\ln   \frac{\dd  \mathbb{Q} }{\dd  \mathbb{P}}(\mH^{g_t}_{0:T}) & =  -\frac{1}{2}\int_0^T \sigma_t^{2} \| h_t(\mH_t^{g_t}) \|^2 \dd t  +   \int_0^T \sigma_t^2 (g_t^\top h_t) (\mH_t^{g_t}) \dd t - \ln p(\vy| \vx_0^{g_t} )  \nonumber\\
&+  \int_0^T \sigma_t h_t^\top (\mH_t^{g_t})\bwd{ \dd \rv{W}}_t,
\end{align}
for the reference process $\mathbb{W}$. The core advantage of VarGrad is that we can choose this reference process. In particular, the choice $g_t=\mathrm{stop\_grad}(h_t)$ gives us a way to detach the trajectories, saving us from having to score all gradients in memory. 

%\begin{align}
%  D_{\mathrm{logvar}}(\mathbb{Q}, \mathbb{P} ; \mathbb{W}) =\mathbb{E}_{\mH^{2 \beta_t \mathrm{stop\_grad}(h_t)}_{0:T} \sim \mathbb{W}} \left[\left(\ln   \frac{\dd  \mathbb{Q} }{\dd  \mathbb{P}}(\mH^{2 \beta_t \mathrm{stop\_grad}(h_t)}_{0:T})  - \mathbb{E}\left[\ln   \frac{\dd  \mathbb{Q} }{\dd  \mathbb{P}}(\mH^{2 \beta_t \mathrm{stop\_grad}(f_t)}_{0:T})\right] \right)^2\right]
%\end{align}

%This allows us to detach the trajectories $\mH^{2 \beta_t \mathrm{stop\_grad}(f_t)}_{0:T}$  saving us from having to store $K * \mathrm{dim}(\vx_0)$ gradients in memory which is challenging for images.

\paragraph{Trajectory Balance} An alternative to the VarGrad loss in Eqn.~\eqref{eq:var_grad_stochastic_control} is the following trajectory balance \cite{malkin2022trajectory,malkin2022gflownets} loss
\begin{align}
    \label{eq:tbloss}
  \mathcal{L}^{\mathbb{W}}_{\mathrm{TB}}(\mathbb{Q}, \mathbb{P} ; k) = \mathbb{E}_{\mH^{g_t}_{0:T} \sim \mathbb{W}} \left[\left(\ln   \frac{\dd  \mathbb{Q} }{\dd  \mathbb{P}}\left(\mH^{g_t}_{0:T}\right)  - k \right)^2\right],
\end{align}
where $k \in \R$ is a learnable parameter and $\mathbb{W}$ is the same reference process as above. We again choose $g_t = \mathrm{stop\_grad}(h_t)$ This loss is also motivated by a valid divergence, see \cite{nusken2021solving}. In difference to the VarGrad loss~\eqref{eq:var_grad_stochastic_control}, the inner expectation is exchanged with $k$, which approximates a running mean. In practice, we optimise $k$ and and $h_t$ at the same time. 

\paragraph{Trajectory Subsampling} We observed that backpropagating gradients for only a random subset of discrete timesteps in the RND can be an effective strategy. This subsampling reduces memory costs at the expense of increased variance and the introduction of a small bias. Nevertheless, the reduced memory usage per example enables larger batch sizes, which can mitigate these effects. Notably, we found that backpropagating for only $20\%$ of the timesteps still achieves comparable performance.


\subsection{Stochastic Optimal Control - Experiments}
\label{sec:stoch_cont_exp}
We provide some initial proof of concept experiments on the MNIST dataset \cite{mnisthandwrittendigit} of handwritten digits. In particular, we make use of a parallel beam Radon transform with $5$ angles as the forward operator and perturb the observations with $10\%$ additive Gaussian noise. All SDEs are discretised using an Euler-Maruyma scheme and the integrals are estimated using simple quadrature rules. We use a non-equidistant time grid according to a square root function, i.e., let $0=t_0 \le \dots \le t_{K-1}=1$ be an equidistant grid of $[0,1]$ for $K$ time points. We then use $t_0^2, t_1^2, \dots, t_{K-1}^2$ as the time points for evaluating the SDEs. This gives us a finer discretisation closer to $t=0$. The unconditional MNIST model is based on the attention U-Net architecture \cite{dhariwal2021diffusion} with about $3$M parameters. The $h$-transform is implemented using the same DEFT parametrisation as in Section~\ref{sec:network_architecture} with about \num{70000} parameters. We compare \texttt{EM-backprop}, i.e., directly backpropagting through the discrete SDE solver (see e.g. \citep{li2020scalable}), with \texttt{VarGrad} and \texttt{TrajectoryBalance} from Section~\ref{app:vargrad}. For \texttt{EM-backprop} we use a batch size of $16$ and use $60$-$80$ time steps. Instead, for \texttt{VarGrad} and \texttt{TrajectoryBalance} we were able to use a batch size of $26$ and $80$-$140$ time steps. For training we used a single GeForce RTX 3090 and the training time took about $1$h. The results are presented in Figure~\ref{fig:stoch_control_loss}, where the loss trajectory of Eqn.~\eqref{eq:stoch} and the mean PSNR of samples is shown. Both \texttt{VarGrad} and \texttt{TrajectoryBalance} are able to minimise the stochastic optimal control objective to the same extend as \texttt{EM-backprop}. Further, in Figure~\ref{fig:trajectory_balance_training} we provide an overview of the \texttt{TrajectoryBalance} training. Here, we observe that $k$ is working as a estimator of the mean RND with a lower variance.

\begin{figure}[t]
    \centering
    \includegraphics[width=0.95\textwidth]{"/Users/shreyaspadhy/Documents/Obsidian Vault/phd-thesis/Chapter5/Figs/Raster/Figs/stochastic_optimal_control_loss.pdf"}
    \caption{Left: Tracking the stochastic optimal control loss \eqref{eq:stoch} for the three methods. Right: Mean PSNR of samples.}
    \label{fig:stoch_control_loss}
\end{figure}

\begin{figure}[t]
    \centering
    \includegraphics[width=0.95\textwidth]{"/Users/shreyaspadhy/Documents/Obsidian Vault/phd-thesis/Chapter5/Figs/Raster/Figs/tb_loss.pdf"}
    \caption{Training using \texttt{TrajectoryBalance}. Left: The trajectory balance objective loss \eqref{eq:tbloss} over training steps. Right: The mean RND and the trained $k$. We see that $k$ follows the mean RND. However, it has a smaller variance.}
    \label{fig:trajectory_balance_training}
\end{figure}


\section{Conclusion}
\label{sec:conclusion}
We presented a unified mathematical framework, based on Doob's $h$-transform, to better understand and classify different conditional diffusion methods.
Under this framework, we proposed DEFT, a novel parameter-efficient conditional fine-tuning method that does not require backpropagation through large pre-trained score networks, resulting in efficient sampling. We evaluated DEFT on several image reconstruction tasks and showed that it reliably outperformed standard methods, both in time, reconstruction quality and perceptual similarity metrics. While DEFT requires additional training on a small dataset of paired measurements, we find that it is still faster than many existing baselines due to being able to use fewer sampling steps during evaluation, and not needing to backpropagate during evaluation.

\vspace{-0.4cm}
\paragraph*{Limitations and future work} The DEFT framework uses
% relies on the availability of 
a (small) fine-tuning dataset, 
% which is 
in contrast to zero-shot conditional sampling approaches, e.g., DPS \cite{chung2022diffusion} or $\Pi$GDM \cite{song2022pseudoinverse}. Fine-tuning on small datasets may have the risk of 
% reproducing and 
overfitting to biases inherent in the data. 
In contrast to zero-shot conditional sampling, DEFT assumes no knowledge of the forward operator. However, the forward operator can be incorporated as an inductive bias within the network architecture to improve performance.
% 
We also proposed a zero-shot approach through the optimal control loss in Section~\ref{sec:online_finetuning}, which only needs a single observation $\m$ to learn the $h$-transform. Though we show promising results scaling this approach to the MNIST dataset in Appendix~\ref{app:amort_cond_training}, the computational burden of simulating the full SDE at each iteration is still high, which might make this optimal control loss infeasible for high-dimensional data. However, there is promising recent work on partial trajectory optimisation \cite{zhang2023diffusion}, which may reduce the computational burden of the stochastic control objective, making it competitive with existing methods.